% VERBATIM TRANSFER ARCHIVE. This is not a standalone paper.
% Do not input every block together: labels retained for provenance may duplicate
% labels kept in the focused note. Move only the block required at its named destination.

% === 01_remove_obsolete_editorial_directives; original lines 1--34 ===
% Destination: Editorial archive only; not mathematical content.
% SPDX-FileCopyrightText: 2026 Will Cook
% SPDX-License-Identifier: CC-BY-4.0
%
% PROVENANCE: superseded historical draft.  The published successor is
% paper/erdos-249-binary-totient-series.tex in the public release repository
% wcook04/plectis-lean-erdos249-257 at 8df0dda280 (release/final-20260905).
% Every label in this file is also defined there; do not export or promote
% from this copy.
%
% Standalone Erdős Problem Note for #249.  This deliberately repeats the
% gateway paper's definitions and status boundary: independent readability is
% more important than avoiding duplication.
%
% Ordering rule.  Most readers stop after two pages, and one of them is a
% moderator who has already rejected a longer version of this work for
% overstating scope.  So: the title names a theorem; the abstract states the
% results and then the boundary; Section 1 is a numbered inventory of every
% claim, each carrying a status tag and a live link into the checked source.
% Everything after Section 1 is mechanism, scope and provenance.
%
% Second rule, learned from that rejection: a barrier is only worth printing
% if its hypothesis class is one somebody would actually assume.  Two
% propositions that appeared here in earlier drafts are demoted in Section 6
% for failing that test, with the reason stated rather than the propositions
% quietly dropped.
%
% Third rule, July 2026.  The status tag on a result must be the status the
% claim registry records for the row that owns its declarations, not a tag
% chosen while writing.  Six of the nine tags in the previous draft disagreed
% with the registry, and two headline results turned out to be owned by no row
% at all.  Those now say so on their face.  A note whose subject is the gap
% between checked propositions and public wording cannot afford that gap in
% its own status column.


% === 02_compact_abstract; original lines 64--126 ===
% Destination: a249_front.tex: the old inventory belongs beside the theorem catalogue, not in an abstract.
\begin{abstract}
\noindent
For
$j\ge0$ and $0\le r<2^j$, write $\ph_{j,r}(n)=\ph(2^jn+r)$.  We prove that
\[
  \{\ph_{0,0},\ph_{1,0}\}
  \cup
  \{\ph_{j,r}:j\ge1,\ 0<r<2^j,\ r\ {\rm odd}\}
\]
is a $\Q$-basis for the span of the $2$-kernel of Euler's totient.  Thus the
dyadic sections through level $e$ have rank exactly $2^e+1$ for every $e\ge1$,
and two elementary identities generate all their rational linear relations.
Independence is read off a finite evaluation matrix that is diagonal modulo a chosen prime; the same argument gives rank $k^e+1$ in every integer base $k\ge2$.  The resulting infinite-dimensionality is also implied by Coons's non-$2$-regularity theorem.

Separately, the least-residue series $A_m=\sum_{n}(\ph(n)\bmod m)\,2^{-n}$
is irrational for every $m\ge3$.  At dyadic resolution $m=2^k$, a
rational-valued map on residue classes of $\ph(n)$ has rational binary value
if and only if it is constant on the even-indexed classes; the irrationality
half is locally proved, the converse is an ordinary identification of the
recurrent support.  Even-indexed means residue classes of $\ph(n)$ modulo
$2^k$, not the odd-residue kernel sections of the basis above.  Neither
statement decides the unreduced series $\Stot$.

A second unconditional theorem determines the monomial rank-one M\"obius--Mersenne quotients $Q(e,Y)$.  The unique minimizer is
$(e,Y)=(1,5)$; every admissible quotient and every nonempty finite positive
mixture exceeds $\Theta_2=\Stot-1/2$ by more than $21/320$.  This determines the monomial quotients $Q(e,Y)$ and their positive finite averages; signed coefficients and coupling before the quotient remain outside the statement.

The closest attack-changing result for the open series is a binary-cyclotomic
anchor theorem.  Clean anchors and unbounded prime-divisor support occur
unconditionally in the actual binary layers, but that support can coexist with
the modular period lock forced by hypothetical rationality.  Thus prime
production is not the missing step: one must produce cofinal phase
discrepancy.  The full-depth amplifier adds that once a single nonintegral seed
exists on a ray, eventual two-syndetic full-depth certificates follow.  It
removes depth as a cost without producing the first seed.

On the size of a displayed denominator the current unrestricted floor is
$q\ge2^{39989}$ and $q>10^{12038}$, from $23438$ certified continued-fraction partial
quotients computed from an $80000$-bit exact rational bracket.  Bit length and decimal digit count are not those exponents, and $10^{12038}<2^{39990}<10^{12039}$, so the two comparisons are not a converted chain.  That is an exact computation and not a Lean theorem.  The Lean-checked Farey exclusion
below, $q>7.96\times10^{34}$, is the formalised instance of a weaker family;
both are functions of how many binary digits were committed rather than
measures of how far the problem has moved.

These advances do not settle Erd\H{o}s Problem~\#249, which asks whether
\[
  \Stot=\sum_{n\ge1}\frac{\ph(n)}{2^n}
\]
is irrational.  For $H_a=\lcm(1,\ldots,2^a)$ and
$R_a=\operatorname{totientTail}(2H_a)-\operatorname{totientTail}(H_a)$, we
prove the exact criterion
\[
  \Stot\notin\Q
  \quad\Longleftrightarrow\quad
  \forall a_0\ \exists a\ge a_0, R_a\notin\Z
\]
The actual-LCM and canonical Mersenne formulas are exact coordinates for that
same unresolved burden, not stronger progress than the anchor/no-go and
amplification theorems above. Two quantitative routes also remain
conditional: one separates the actual LCM orbit from the integers; the other
controls four terms in a first-harmonic decomposition. Their cofinal hypotheses
are unproved. The exact open boundary is therefore arithmetic control of the
actual tail, not the finite-rank structure of the totient's dyadic sections.
\end{abstract}

% === 03_remove_repeated_result_panel; original lines 128--140 ===
% Destination: Editorial archive only; all mathematical assertions retained in catalogue.
\begin{paperresult}{What is proved here}
\textbf{Unconditional structure.} The displayed family is a rational basis
for the $2$-kernel of Euler's totient, so the sections through level $e$
have rank exactly $2^e+1$.  The positive rank-one
M\"obius--Mersenne cone also has an exact unique minimizer and a uniform gap.
\textbf{Anchor and no-go advance.} The actual binary-cyclotomic layers have
clean anchors and unbounded prime-divisor support, but support alone can remain
period-locked; a separate amplifier removes full depth as a cost once a seed
exists.
\textbf{Series reduction.} Irrationality is equivalent to nonintegrality of
the actual LCM-tail differences cofinally.  \textbf{Open boundary.} The basis
theorem and cone obstruction do not supply even one required nonintegral seed.
\end{paperresult}

% === 04_lead_with_existing_basis_and_full_proof; original lines 142--245 ===
% Destination: Definitions reused below. Prior introduction/status history: a249_front.tex after prop:rank; reduction examples retained in new opening.
\section{Introduction and main results}\label{sec:results}

Erd\H{o}s Problem~\#249 asks whether
\[
  \Stot=\sum_{n\ge1}\frac{\ph(n)}{2^n}
\]
is irrational; see Erd\H{o}s and Graham~\cite[p.~61]{erdosgraham1980} and
Erd\H{o}s~\cite[p.~102]{erdos1988}.  Bloom's current catalogue record
reproduces this question and labels it open, while explicitly warning that
the status is the website owner's present assessment and may omit relevant
literature~\cite{erdosproblems}.  We therefore use the catalogue for
numbering and current reported status only; the two original sources carry
the problem statement.

The principal unconditional result is structural: the two zero-residue
sections $\ph_{0,0}$ and $\ph_{1,0}$, together with the odd-residue sections,
form a rational basis, and the two reduction identities below generate every
rational relation among the sections.  It gives exact finite-level rank but
no rationality-to-finite-rank bridge.  The monomial rank-one M\"obius--Mersenne quotients $Q(e,Y)$ are also determined sharply: one unique minimizer, the uniform $21/320$ floor, the $1/16$ versus $1/15$ boundary, and closure under arbitrary finite positive averages of those quotients.

The closest results concerning the open series are exact reformulations and
conditional criteria.  The actual power-two LCM orbit gives a cofinal
nonintegrality equivalence; the canonical Mersenne formulation removes the
remote basepoint; and the full-depth theorem shows that one seed gives eventual
$2$-syndeticity on its ray.  The independent first-harmonic argument turns four
explicit fibre bounds into a $9X/10$ window estimate, but irrationality requires
their unproved cofinal form.  The example $38=2\cdot19$ rules out the simplest
global-isolation argument.  The remaining inventory records finite Farey and
diagonal certificates, Lambert and gcd-moment identities, and explicit limits
of fixed-precision methods.  None supplies the required cofinal input.

\statusboundary

Each \emph{Checked} line below links the relevant declarations at
commit~\texttt{\commitshort}.  The bracketed tags reproduce the public status
recorded by the claim registry; where no registry row owns the cited
declarations, the text says so rather than assigning a neighbouring status.

Throughout $\N=\{0,1,2,\ldots\}$ and $\ph(0)=0$.  Thus
$\ph_{j,r}(n)=\ph(2^jn+r)$, for $j\ge0$ and $0\le r<2^j$, is the
$(j,r)$ \emph{dyadic section} of Euler's totient, a vector in $\Q^{\N}$, and
$j$ is its \emph{level}; the family of all of them is the $2$-kernel
$\Ker_2(\ph)$ in the sense of Allouche and Shallit~\cite{allouche-shallit}.
Two elementary identities relate them: the \emph{reductions}
$\ph_{j+1,0}=2^{j}\ph_{1,0}$
(\mword{Erdos249257/TotientMahlerDefect.lean}{160}{totientKernel_zero_residue}{checked})
and $\ph_{j,\,2^{t+1}s}=2^{t}\ph_{j-t-1,\,s}$ for $s$ odd with $2^{t+1}s<2^j$
(\mword{Erdos249257/TotientMahlerDefect.lean}{169}{totientKernel_even_residue_reduce}{checked}).
At level $3$, for instance, they give $\ph_{3,0}=4\ph_{1,0}$,
$\ph_{3,2}=\ph_{2,1}$, $\ph_{3,4}=2\ph_{1,1}$ and $\ph_{3,6}=\ph_{2,3}$, while
the four odd residues $r=1,3,5,7$ are reduced by neither; numerically
$\ph_{3,4}(1)=\ph(12)=4=2\ph(3)=2\ph_{1,1}(1)$.  So beyond $\ph_{0,0}$ and
$\ph_{1,0}$ each level contributes only its odd residues, and the count
through level $e$ is $2+\sum_{j=1}^{e}2^{j-1}=2^e+1$; \ref{res:rank} says that
count is the exact dimension.  Independence is the finite evaluation matrix of Section~\ref{sec:rank}, diagonal modulo a prime.

Two objects built from $\Stot$ recur below.  The \emph{scaled tail}
$R_N=\sum_{m\ge1}\ph(N+m)/2^m$ is $2^N$ times the part of the series beyond
index $N$; since $2^N\Stot=\Phi_N+R_N$ with
$\Phi_N=\sum_{1\le n\le N}\ph(n)2^{N-n}$ an integer
(\mword{Erdos249257/TotientTailPeriodKiller.lean}{150}{two_pow_mul_totient_series_eq}{checked}),
$R_N$ and $2^N\Stot$ differ by an integer.  For $h,N,L\in\N$, put
\[
  D_{h,N,L}
  =
  \sum_{j=0}^{L-1}
  \bigl(\ph(N+h+1+j)-\ph(N+1+j)\bigr)2^{L-1-j},
\]
and let $\rho_{h,N,L}\in\{0,\dots,2^L-1\}$ be its residue modulo $2^L$.
A \emph{finite tail-difference certificate} at $(h,N,L)$ is the pair of
strict inequalities
\[
  N+h+L+2<\rho_{h,N,L}<2^L-(N+h+L+2)
\]
(\mword{Erdos249257/TotientTailPeriodKiller.lean}{72}{certifiedKill}{definition}).
It exists to make a statement about an infinite tail decidable by a finite
computation, and it is faithful: some depth $L$ certifies $(h,N)$ if and only
if $R_{N+h}-R_N\notin\Z$
(\mword{Erdos249257/LcmConeFlatness.lean}{316}{exists_certifiedKill_iff_tail_diff_notMem_int}{checked}).
For example,
\[
  D_{1,12,16}=-143140,\qquad
  \rho_{1,12,16}=53468,\qquad N+h+L+2=31,
\]
so $(1,12,16)$ is a certificate.  More generally, the test succeeds at
$(h,12,16)$ for every $h$ with $1\le h\le8$
(\mword{Erdos249257/TotientTailPeriodKiller.lean}{404}{certifiedKill_all_small}{checked}),
so none of $R_{13}-R_{12},\dots,R_{20}-R_{12}$ is an integer
(\mword{Erdos249257/TotientTailPeriodKiller.lean}{407}{tail_diff_not_int_small}{checked}).

A third recurring object is an \emph{integral scaled-tail sequence with
multiplier $v$}.  For a coefficient sequence $c:\N\to\N$ and an integer
$v\ge1$, this
is an integer sequence $u$ satisfying
\[
  u(N+1)=2u(N)-v\,c(N+1)\quad\hbox{for every }N,
  \qquad
  \frac{u(N)}{2^N}\longrightarrow0
\]
(\mword{Erdos249257/GenericTailOrbitRigidity.lean}{67}{IsTemperedBinaryOrbit}{definition});
for $c=\ph$ such a pair $(v,u)$ with $v\ge1$ exists exactly when $\Stot$ is
rational (Section~\ref{sec:carry-rank}).



% === 05_transfer_result_inventory; original lines 245--730 ===
% Destination: a249_front.tex after prop:rank for structural results; a249_p1b.tex prop:TE-05 for canonical/full-depth; route-specific destinations in LONG_RECORD_MAP.md.
{\small
\begin{enumerate}[label=\textbf{R\arabic*.},ref=R\arabic*,leftmargin=2.4em,
  itemsep=0.2em,topsep=0.25em]

\item\label{res:canonicalmersenne} \stag{formalised here}
\textbf{Canonical Mersenne arithmetic frontier.}  Let
\[
 B_{H,N}=\sum_{j=0}^{H-1}\ph(N+1+j)2^{H-1-j},\qquad
 \rho_{H,N,M}=(-B_{H,N})\bmod M.
\]
For every $c,v\in\N$ with $v>0$ and $\gcd(2,v)=1$, the statement that there
are $H,M\in\N$ with
\[
 H>0,\qquad \ph(v)\mid H,\qquad vM=2^H-1,
 \qquad c+H+1<\rho_{H,c,M}<M-(c+H+1)
\]
is equivalent, after quantifying over $c,v$, to irrationality of $\Stot$.
	The basepoint is exactly $c$, the dyadic valuation of a prospective rational
	denominator; the older centred-residue condition searched for a remote
	$N\ge\max(c,N_0)$.  The canonical statement implies that older condition by
feeding $\max(c,N_0)$ into its $c$-slot.  Moreover, whenever $M\mid2^H-1$,
the finite residue obeys the exact recurrence
\[
 \rho_{H,N+1,M}
 =\bigl(2\rho_{H,N,M}-(\ph(N+H+1)-\ph(N+1))\bigr)\bmod M.
\]
This removes analytic tails, carry trajectories, and remote-basepoint search
from the equivalent residual, but it does not prove the required central gap.
{\footnotesize\emph{Checked:}
\lword{Erdos249/CyclotomicAnchoredKill.lean}{2322}{fullMersenneBlockResidue_succ}{residue recurrence},
\lword{Erdos249/CyclotomicAnchoredKill.lean}{3132}{fullMersenneCenteredResidueGapSupply_of_canonicalBasepoint}{dominance over the remote condition},
\lword{Erdos249/CyclotomicAnchoredKill.lean}{3165}{fullMersenneCanonicalBasepointResidueGapSupply_iff_irrational}{exact equivalence}.}

\item\label{res:fulldepth} \stag{formalised here}
\textbf{Full-depth ray amplification and exactness.}  Fix $d>0$ and $N$.
If $(d,N,L)$ is a finite tail-difference certificate at any depth $L$, then
there is $T>0$ such that every pair $\{t,t+1\}$ with $t\ge T$ contains a
multiplier $m$ for which $(md,N,md)$ is a certificate.  Thus one seed forces
an eventually $2$-syndetic family whose depth equals its period.  Pointwise,
\[
 \exists t>0:\ (td,N,td)\text{ is certified}
 \quad\Longleftrightarrow\quad R_{N+d}-R_N\notin\Z.
\]
Consequently the predicate demanding such a full-depth certificate for every
$d>0,N$ is equivalent to irrationality of $\Stot$.  Its cofinal-in-$N$
	version is equivalent both to the older arbitrary-depth period-multiple
	condition and to irrationality.  Thus requiring full depth loses no strength;
it does not produce the first non-integral shift or seed.
{\footnotesize\emph{Checked:}
\lword{Erdos249/FullDepthRayAmplifier.lean}{364}{eventually_twoSyndetic_fullDepthKillMultipliers_of_seed}{eventual two-syndeticity},
\lword{Erdos249/FullDepthRayAmplifier.lean}{388}{exists_fullDepthKill_on_ray_iff_shift_notMem_int}{pointwise equivalence},
\lword{Erdos249/FullDepthRayAmplifier.lean}{406}{apFullDepthEscape_iff_irrational}{global equivalence},
\lword{Erdos249/FullDepthRayAmplifier.lean}{421}{cofinalFullDepthKillSupply_iff_periodMultipleKillSupply}{cofinal depth deletion},
\lword{Erdos249/FullDepthRayAmplifier.lean}{438}{cofinalFullDepthKillSupply_iff_irrational}{cofinal endpoint}.}

\item\label{res:basis} \stag{unconditional progress}
\textbf{Basis theorem.}  The family
$\mathcal B=\{\ph_{0,0},\ph_{1,0}\}\cup\{\ph_{j,r}:j\ge1,\ r\ \text{odd},\
0<r<2^j\}$ is linearly independent over $\Q$ and spans the same $\Q$-subspace
of $\Q^{\N}$ as $\Ker_2(\ph)$; so $\mathcal B$ is a basis of that span, which
is therefore infinite-dimensional.  Since the reductions make every remaining
section a rational multiple of a member of $\mathcal B$, the $\Q$-linear
relations among dyadic sections of $\ph$ are exactly those the reductions
generate, a one-line consequence of the checked statements, not a further
one.
{\footnotesize\emph{Checked:}
\mword{Erdos249257/TotientMahlerDefect.lean}{1265}{linearIndependent_oddCoreTotientKernelFamily}{independence},
\mword{Erdos249257/TotientMahlerDefect.lean}{1380}{span_range_fullTotientKernel_eq_span_range_oddCore}{span equality},
\mword{Erdos249257/TotientMahlerDefect.lean}{1392}{totientDyadicSectionBasis}{basis}.
\emph{Registry:} \texttt{dyadic\_totient\_basis}.}

\item\label{res:rank} \stag{unconditional progress}
\textbf{Exact finite-level rank.}  For every $e\ge0$ the $2^e+1$ sections
$\ph_{0,0}$, $\ph_{1,0}$ and $\ph_{j,r}$ with $1\le j\le e$, $r$ odd,
$0<r<2^j$ are linearly independent over $\Q$, so their span has dimension
exactly $2^e+1=2+\sum_{j=1}^{e}2^{j-1}$: exact, not an estimate, because the
even residues are already dependent by the reductions.  Read instead as a
truncation of the kernel, $V_e=\operatorname{span}_{\Q}\{\ph_{j,r}:0\le j\le
e\}$ has $\dim V_e=2^e+1$ for every $e\ge1$, and $\dim V_0=1$: the counted
family carries $\ph_{1,0}$, which is a level-one section, so the two readings
agree from $e=1$ on and differ only at $e=0$.  At $e=1$ both readings say that
$\ph(n)$, $\ph(2n)$ and $\ph(2n+1)$ are linearly independent over $\Q$.  The
independence half of
\ref{res:basis} follows by finite character; the span equality needs the
reductions as well.
{\footnotesize\emph{Checked:}
\mword{Erdos249257/TotientMahlerDefect.lean}{935}{linearIndependent_canonicalTotientKernelFamily}{independence},
\mword{Erdos249257/TotientMahlerDefect.lean}{989}{finrank_canonicalTotientKernel_eq}{dimension},
\mword{Erdos249257/TotientMahlerDefect.lean}{1084}{finrank_totientKernelThroughLevelFamily_eq}{the exact finite-truncation rank};
the latter is the complete truncation wrapper, with its explicit hypothesis $e\ge1$,
while the canonical-family rank is the intermediary because even residues reduce to
the odd-core family.
\mword{Erdos249257/TotientMahlerDefect.lean}{1145}{not_finiteDimensional_span_fullTotientKernel}{infinite-dimensionality directly}.}

\item\label{res:residueseries} \stag{unconditional progress}
\textbf{Least-residue series and even-indexed dyadic classification.}
For every $m\ge3$ the series $A_m=\sum(\ph(n)\bmod m)\,2^{-n}$ is
irrational.  At $m=2^k$ ($k\ge1$), a rational-valued letter map on residues
of $\ph$ has rational binary value if and only if it is constant on the
even-indexed residue classes.  The forward half is Lean-checked; the converse
is the ordinary identification that every even residue is recurrent at dyadic
moduli.  Even-indexed means classes of $\ph(n)$ modulo $2^k$; it is not the
odd-residue indexing of \ref{res:basis}.  An ordinary all-modulus
classification (rational iff constant on the recurrent set $\mathcal R_m$)
is recorded separately and does not imply irrationality of $\Stot$.
{\footnotesize\emph{Checked:}
\lword{Erdos249/ResidueClassTotientSeries.lean}{518}{irrational_totientObservable_of_isolated_residue}{pulse consumer without a unit centre},
\lword{Erdos249/ResidueClassTotientSeries.lean}{553}{irrational_totientObservable}{unit-centre special value},
\lword{Erdos249/ResidueClassTotientSeries.lean}{569}{fixed_resolution_observable_irrational}{dyadic forward half},
\lword{Erdos249/ResidueClassTotientSeries.lean}{595}{residue_series_irrational}{$A_m$ for $m\ge3$}.
\emph{Registry:} \texttt{residue\_class\_totient\_series\_irrational}.}

\item\label{res:rankonefloor} \stag{unconditional progress}
\textbf{Sharp positive rank-one and direct-sum floor.}  Put
\[
 \Theta_2=\sum_{d\ge1}\frac{\mu(d)}{(2^d-1)^2}=\Stot-\frac12,
 \qquad
 Q(e,Y)=
 \frac{\bigl(\sum_{d=1}^{Y}\mu(d)/(2^d-1)^{e+2}\bigr)^2}
      {\sum_{d=1}^{Y}\mu(d)/(2^d-1)^{2e+2}}.
\]
The identity $\Stot=\sum_{d\ge1}\mu(d)/(2^d-1)^2+1/2$ is Steve Fan's, posted on
the problem's discussion page on 16 May 2026~\cite{fan2026totient}; the object
$\Theta_2$ is his and only the extremal calculation below is ours.
For every $e\ge1$ and $Y\ge4$, the quotient $Q(e,Y)$ is uniquely minimized
at $(1,5)$ and satisfies $Q(e,Y)-\Theta_2>21/320>1/16$.  The stronger
unit-fraction floor $1/15$ fails already at $(1,5)$.  Every nonempty finite
positive weighted average of admissible $Q(e,Y)$ has the same strict
$21/320$ separation.  In rational-linear-form language, if
$Q(e,Y)=p/q$ with $p\in\Z$ and $q\ge1$, then
$|q\Theta_2-p|>21q/320$.  Hence neither adding atoms nor taking positive direct
sums can repair this rank-one approximation mechanism; signed cancellation
and genuinely coupled higher-rank kernels remain outside the theorem.
{\footnotesize\emph{Checked:}
\lword{Erdos249/RankOneSharpFloor.lean}{534}{rankOneSubrankQuotient_eq_one_five_iff}{unique minimizer},
\lword{Erdos249/RankOneSharpFloor.lean}{624}{rankOneSubrankQuotient_sub_theta_two_gt_twentyOne_div_threeTwenty}{uniform $21/320$ floor},
\lword{Erdos249/RankOneSharpFloor.lean}{634}{rankOneSubrankQuotient_sub_theta_two_gt_one_div_sixteen}{$1/16$ floor},
\lword{Erdos249/RankOneSharpFloor.lean}{645}{not_forall_rankOneSubrankQuotient_sub_theta_two_gt_one_div_fifteen}{$1/15$ failure},
\lword{Erdos249/RankOneSharpFloor.lean}{655}{positive_direct_sum_sub_theta_two_gt_twentyOne_div_threeTwenty}{positive direct sums},
\lword{Erdos249/RankOneSharpFloor.lean}{693}{primitive_form_abs_gt_twentyOne_div_threeTwenty}{rational linear-form obstruction}.}

\item\label{res:index} \stag{proved here}
\textbf{Canonical all-base index and unique coordinates.}  For every integer
$k\ge2$ and level $e$, the two zero-residue channels are followed at level
$j+1$ by the unique coordinates
$r=kq+(d+1)$ with $q<k^j$ and $d<k-1$.  Equivalently, every integer
$1\le r<k^{j+1}$ with $k\nmid r$ has one and only one such representation.
The level block therefore has $k^{j+1}-k^j$ coordinates, and the complete
finite index has exactly
$2+\sum_{j<e}k^j(k-1)=k^e+1$ elements.  This is the concrete quotient and
nonzero-last-digit mechanism behind the all-base rank statement: it counts
the canonical candidates and proves the reduction coordinates, but it does
not prove their linear independence.
{\footnotesize\emph{Evidence:}
\mword{Erdos249257/AllBaseTotientKernel.lean}{852}{card_allBaseCanonicalIndex}{index count},
\mword{Erdos249257/AllBaseTotientKernel.lean}{644}{linearIndependent_totientAffineForms}{independence}.
The count is a live Lean theorem for $e\ge1$; independence is unconditional in the same module.  Martin remains prior art, not a Lean dependency.}

\item\label{res:denominator} \stag{formalised here}
\textbf{Denominator exclusion, sharp for its window.}  If $\Stot=a/q$ with
$a\in\Z$ and $q\ge1$ then
$q>79\,639\,646\,646\,701\,375\,323\,355\,774\,875\,831\,053\approx7.96\times10^{34}$.
The constant is optimal for the window that produces it, the window here being
the pair $(N,K)=(1,240)$, meaning $240$ committed binary digits of the series
shifted by one, in the sense that the next integer up is the exact first
denominator at which that window's certificate fails.  The method is the
classical Farey mediant argument (Section~\ref{sec:denominator}).
{\footnotesize\emph{Checked:}
\mword{Erdos249257/CertificateKernel.lean}{18371}{tsum_totient_div_pow_two_ne_int_div_of_den_le_79639646646701375323355774875831053}{exclusion},
\mword{Erdos249257/CertificateKernel.lean}{18384}{tsum_totient_div_pow_two_ne_ratCast_of_den_le_79639646646701375323355774875831053}{reduced-denominator form},
\mword{Erdos249257/GapFareyBound.lean}{225}{gap_check_window_1_240_first_failure}{first failure}.
\emph{Registry:} the row owns the reduced-denominator form only.}

\item\label{res:carryrank} \stag{formalised here}
\textbf{Rank lower bound for a rational scaled-tail sequence.}  If $\Stot$
were rational then some integral scaled-tail sequence $u$ for the coefficients
$\ph$, with positive multiplier, would have for \emph{every} $e$ dyadic
sections through level $e$, formed from $u$ exactly as the $\ph_{j,r}$ are
formed from $\ph$, spanning a space of dimension at least $2^e-1$.  The
same lower bound holds for every positive-multiplier integral scaled-tail
sequence.  This statement is specific to that
binary scaled-tail representation; it does not constrain arguments formulated
only through the Lambert weight, the coprimality identity, or tail differences.
It is not an irrationality criterion, because no finite-rank upper bound is
proved for the scaled-tail sequences supplied by rationality.
{\footnotesize\emph{Checked:}
\mword{Erdos249257/TotientCarryKernelRigidity.lean}{300}{not_irrational_totientSeries_implies_unbounded_carryRank_unconditional}{the barrier},
\mword{Erdos249257/TotientCarryKernelRigidity.lean}{211}{finrank_canonicalCarryKernel_ge_of_linearIndependent}{rank transport from \ref{res:rank}}.}

The orbit-level form sharpens this rank barrier without closing it: under the
same non-irrationality hypothesis, the same carry also has uniformly eventual
dyadic-section periodicity modulo $v$, as recorded by
\mword{Erdos249257/TotientTailCarryPeriod.lean}{224}{not_irrational_totientSeries_implies_mod_period_and_unbounded_rank}{the modular-period/rank frontier}.
This does not turn quotient periodicity into a finite-$\mathbb Q$-rank bound.
Indeed, if $p\mid v$, the totient forcing disappears modulo $p$ and the carry
is the homogeneous geometric orbit $u_N\equiv2^N u_0$ modulo $p$
(\mword{Erdos249257/TotientTailCarryPeriod.lean}{126}{totient_forcing_vanishes_mod_of_dvd_multiplier}{forcing vanishing};
\mword{Erdos249257/TotientTailCarryPeriod.lean}{140}{totient_carry_modEq_geometric_of_dvd_multiplier}{geometric reduction}).
This is the natural obstruction to extracting the missing finite-rank upper
bound, so the irrationality of $\Stot$ remains open.

\item\label{res:weight} \stag{checked; no registry row}
\textbf{The \#249 weight.}  $\Stot=\sum_{d\ge1}\Aw(d)/(2^d-1)$ with
$\Aw=\ph*\mu$, the Dirichlet convolution
$\Aw(n)=\sum_{d\mid n}\ph(d)\mu(n/d)$; here $\Aw\ge0$, $\Aw(p)=p-2$,
$\Aw(p^k)=\ph(p^k)-\ph(p^{k-1})$, and $\Aw$ is unbounded.  Its values at
$n=1,\dots,10$ are $1,0,1,1,3,0,5,2,4,0$.
{\footnotesize\emph{Checked:}
\mword{Erdos249257/CertificateKernel.lean}{18445}{tsum_primWeight_div_two_pow_sub_one_eq_totient_series}{identity},
\mword{Erdos249257/MersenneLambertLadder.lean}{370}{primWeight_nonneg}{sign},
\mword{Erdos249257/MersenneLambertLadder.lean}{297}{primWeight_apply_prime}{primes},
\mword{Erdos249257/MersenneLambertLadder.lean}{272}{primWeight_apply_prime_pow}{prime powers},
\mword{Erdos249257/MersenneLambertLadder.lean}{321}{primWeight_not_bounded}{unboundedness}.
\emph{Registry:} the Lambert-ladder row owns the neighbouring rungs, not these
five declarations; see the formal source notes below.}

\item\label{res:gcdmoment} \stag{proved here}
\textbf{Squared Lambert layers expose a first gcd moment.}  The two exact
identities
\[
  \sum_{d\ge1}\frac{1}{(2^d-1)^2}
    =\sum_{n\ge1}(\sigma(n)-\tau(n))2^{-n},
  \qquad
  \sum_{d\ge1}\frac{\ph(d)}{(2^d-1)^2}
    =\sum_{n\ge1}(\mathsf P(n)-n)2^{-n},
\]
where $\mathsf P(n)=\sum_{e\mid n}\ph(e)(n/e)$, are checked in the
source.  For the independent fair-coin geometric waiting times used below,
the second left-hand side is $\mathbb E[\gcd(X,Y)]$: the factor
$(2^d-1)^{-2}$ is the mass of the event $d\mid X$ and $d\mid Y$, and
$\sum_{d\mid g}\ph(d)=g$.  This exact first-moment identity does not imply
irrationality of $\Stot$; no novelty claim is made for the
divisor algebra.
{\footnotesize\emph{Checked:}
\mword{Erdos249257/GcdMomentCalculus.lean}{216}{tsum_one_div_mersenne_sq_eq_sigma_sub_tau_series}{divisor-count layer},
\mword{Erdos249257/GcdMomentCalculus.lean}{235}{tsum_totient_div_mersenne_sq_eq_gcd_moment_series}{first gcd moment}.}

\item\label{res:squareenclosure} \stag{proved here}
\textbf{A sharp squared-Mersenne diagonal enclosure.}  For the diagonal
tail difference \(R_{2H}-R_H\), let
\[
 C_{H,D}=Z_H+2^H(2^H-1)\left(\frac12+
   \sum_{1\le d\le D}\frac{\mu(d)}{(2^d-1)^2}\right),
 \qquad Z_H=\Phi_H-\Phi_{2H}.
\]
The exact decomposition is
\[
 R_{2H}-R_H-C_{H,D}
 =2^H(2^H-1)\sum_{d>D}\frac{\mu(d)}{(2^d-1)^2},
\]
and the signed remainder has the sharp elementary bound
\[
 \left|R_{2H}-R_H-C_{H,D}\right|
 \le \frac{4\,2^H(2^H-1)}{3(2^{D+1}-1)^2}.
\]
Consequently, if this radius is smaller than the distance from \(C_{H,D}\)
to every integer, the diagonal cannot be integral.  The centre is rational
and its denominator is explicitly reduced by the gcd with \(2^H(2^H-1)\);
the separation hypothesis is the remaining arithmetic input.  This finite
criterion gives neither an unbounded separation family nor irrationality of
\(\Stot\).
{\footnotesize\emph{Checked:}
\mword{Erdos249257/SquaredMersenneDiagonalEnclosure.lean}{138}{scaleDiagonalTailDifference_sub_lambertProjectedDiagonal}{exact centre decomposition},
\mword{Erdos249257/SquaredMersenneDiagonalEnclosure.lean}{408}{abs_scaleDiagonalTailDifference_sub_lambertProjectedDiagonal_le}{sharp enclosure},
\mword{Erdos249257/SquaredMersenneDiagonalEnclosure.lean}{426}{scaleFullTarget_miss_of_lambert_projected_separation}{separation criterion};
\mword{Erdos249257/FullTargetPrimeAdjunctionNoGo.lean}{139}{scaleFullTargetHit_iff_integral}{target equivalence}.}

\item\label{res:freshprimedeficit} \stag{proved here}
\textbf{Fresh-prime support is an explicit endpoint deficit.}  For \(H>0\),
put \(g_H(s)=\gcd(\operatorname{rad}(H),s)\) and
\[
 M_H(k,s)=\frac{kH+s}{g_H(s)}\ph(g_H(s)),
 \qquad F_H(k,s)=M_H(k,s)-\ph(kH+s).
\]
The old-prime mass \(M_H(k,s)\) dominates the actual totient, so
\(F_H(k,s)\ge0\).  If \(\Delta_H(s)=\ph(2H+s)-\ph(H+s)\), then
\[
 \Delta_H(s)=\Delta_H^{\rm old}(s)+F_H(1,s)-F_H(2,s),
\]
and the literal foreign channel is exactly the same lower-minus-upper
deficit.  The second-difference and five-point curvature operators preserve
this decomposition; the latter is bounded below by the old-prime branch
minus
\[
 8F_H(1,s)+4F_H(2,s+1)+2F_H(2,s+2)+F_H(2,s+3)+F_H(2,s+4).
\]
This isolates the hard issue: new prime support has a known sign, but no
cofinal bound on its adverse placement is proved.  In particular, this is a
finite curvature-loss budget, not an irrationality proof.
{\footnotesize\emph{Checked:}
\mword{Erdos249257/FreshPrimeDeficitDecomposition.lean}{91}{endpointFreshDeficit_nonneg}{endpoint nonnegativity},
\mword{Erdos249257/FreshPrimeDeficitDecomposition.lean}{184}{finiteForeignChannelIncrement_eq_endpointFreshDeficit_sub}{foreign-channel identity},
\mword{Erdos249257/FreshPrimeDeficitDecomposition.lean}{269}{oldFivePoint_sub_secondBranchAdverseDeficit_le_actual}{five-point loss bound}.}

\item\label{res:squarecrt} \stag{proved here}
\textbf{Finite square-CRT correction suppression.}  Let \(E\) be a finite
family of distinct primes \(p_i\), with anchors \(A_i\), and let \(J<p_i\) for
every \(i\in E\).  A common square-CRT base \(n\) can be chosen with
\[
 n<\left(\prod_{i\in E}p_i^2\right)(\sup_{i\in E}A_i+2),
 \qquad n=A_i+p_i^2t_i
\]
for suitable \(t_i\).  Therefore every \(1\le h\le J\) obeys
\[
 \ph(n+p_i h-A_i)=(p_i-1)\ph(p_i t_i+h).
\]
The congruence removes the \(p_i\)-divisibility correction on the whole
finite horizon, simultaneously across the family.  Checked regression
witnesses show why this is not already a separation theorem: one clean block
has both displayed coefficients zero (at base \(52\)), while another clean
block has a displayed coefficient \(-4\) (at base \(27\)).  Thus the CRT
mechanism controls a finite local error term but supplies neither a cofinal
gap nor an irrationality criterion.
{\footnotesize\emph{Checked:}
\mword{Erdos249257/SquareCRTCube.lean}{297}{exists_squareCRT_clean_totient_family}{simultaneous finite family},
\mword{Erdos249257/SquareCRTCube.lean}{327}{exists_squareCRT_clean_horizon}{finite horizon},
\mword{Erdos249257/SquareCRTCube.lean}{459}{squareCRT_clean_block_can_vanish}{vanishing block},
\mword{Erdos249257/SquareCRTCube.lean}{477}{squareCRT_clean_block_can_be_nonzero}{nonzero block}.}

\item\label{res:slackcriterion} \stag{conditional reduction}
\textbf{Canonical strict-jump slack reduction.}  Put
\[
 \begin{aligned}
 m_t&=\operatorname{canonicalAdjacentSuffixDepth}(t),\\
 d_t&=\operatorname{diagonalAdjacentSuffixResidue}(t,0,m_t),
 \end{aligned}
\]
and define the integer slack
\[
 \lambda_t=\min\!\left(d_t-2^{m_t-5},
       (2^{m_t}-2^{m_t-5})-d_t\right).
\]
The two-sided central-band test at $t$ is equivalent to
$\lambda_t\ge0$.  Moreover, the existence, beyond every threshold, of a
strict LCM jump $H_t<H_{t+1}$ with $\lambda_t\ge0$ is equivalent to the
corresponding central-band condition; either cofinal condition implies that
$\Stot$ is irrational.  This exposes one integer observable for exact
computation and recurrence arguments, but the cofinal sign assertion itself
remains open.
{\footnotesize\emph{Checked:}
\mword{Erdos249257/DiagonalFreshLossBridge.lean}{2667}{canonicalAdjacentSuffixCentral_iff_slack_nonneg}{centrality equivalence},
\mword{Erdos249257/DiagonalFreshLossBridge.lean}{2757}{canonicalAdjacentSuffixJumpCentralSupply_iff_slackSupply}{equivalence of the two conditions},
\mword{Erdos249257/DiagonalFreshLossBridge.lean}{2932}{irrational_totientSeries_of_canonicalAdjacentSuffixJumpSlackSupply}{conditional irrationality theorem}.}

\item\label{res:foreignresidue} \stag{conditional reduction}
\textbf{Finite foreign-residue projection.}  For $H,D\in\N$, split the
finite Möbius-residue diagonal into the channels whose indices divide $H$ and
those that do not.  The checked identity is
\[
 \begin{aligned}
 \operatorname{finiteResidueDiagonal}(H,D)
   &=\operatorname{projectedForeignDefect}(H,D)\\
   &\quad+\operatorname{projectedDivisorChannels}(H,D).
 \end{aligned}
\]
When $H\le D$, the divisor term is exactly the explicit radical Möbius
shadow.  If $2H\le D$, every finite omitted window is bounded by
\[
 B_{H,D}=\operatorname{diagonalCoefficient}(H)\left(
       \frac{2}{2^D}+\frac{4}{3\,4^D}\right).
\]
Thus a controlled comparison of the limiting foreign defect with $B_{H,D}$,
together with separation of the projected state from every integer, forces a
full-target miss.  The analytic comparison and an unbounded supply of
separated projections remain unproved; this is a conditional reduction, not
an irrationality theorem.
{\footnotesize\emph{Checked:}
\mword{Erdos249257/ActualForeignResidueProjection.lean}{276}{abs_foreignTailWindow_le_foreignComplementBound}{finite complement bound},
\mword{Erdos249257/ActualForeignResidueProjection.lean}{308}{finiteResidueDiagonal_eq_projectedForeign_add_divisor}{exact channel split},
\mword{Erdos249257/ActualForeignResidueProjection.lean}{414}{scaleFullTarget_miss_of_projected_separation}{conditional target-miss theorem}.}

\item\label{res:actualorbit} \stag{exact reformulation}
\textbf{Actual LCM orbit frontier.}  Let
$H_a=\lcm(1,\ldots,2^a)$ and
$R_a=\operatorname{totientTail}(2H_a)-\operatorname{totientTail}(H_a)$.
Then $\Stot$ is irrational if and only if, beyond every threshold $a_0$,
some $a\ge a_0$ has $R_a\notin\Z$.  This is an exact cofinal
nonintegrality criterion for the actual LCM diagonal, not a quantitative
separation estimate and not a proof that such indices occur.
{\footnotesize\emph{Checked:}
\mword{Erdos249257/TotientActualLcmOrbitNonintegrality.lean}{37}{irrational_totientSeries_iff_actualLcmOrbitNonintegralitySupply}{exact frontier},
\mword{Erdos249257/TotientActualLcmOrbitNonintegrality.lean}{53}{irrational_totientSeries_of_actualLcmOrbitNonintegralitySupply}{forward implication}.}

\item\label{res:actualsign} \stag{unconditional progress}
\textbf{Positive LCM corridor and top-edge boundary.}  For $a\ge8$ and
$J+(a+6)<2\cdot2^a$, the true translated orbit difference
$\operatorname{totientTail}(2H_a+J)-\operatorname{totientTail}(H_a+J)$ is
strictly positive.  If it is represented by an integer and the dyadic
modulus has room for the directed endpoint strip, the carry survivor is
strictly negative and the discrepancy residue is exactly the top-edge
representative $2^K-\text{carry}_K$, lying strictly between the strip's lower
edge and $2^K$.  Thus the positive corridor eliminates the nonnegative true
survivor but does not itself produce a central-band contradiction; an
independent arithmetic exclusion of the top boundary is still required.
{\footnotesize\emph{Checked:}
\mword{Erdos249257/TotientActualLcmOrbitSign.lean}{39}{actualLcmTailDiff_shift_pos}{positive corridor},
\mword{Erdos249257/TotientActualLcmOrbitSign.lean}{172}{actualLcm_trueEndpointSurvivor_neg}{negative survivor},
\mword{Erdos249257/TotientActualLcmOrbitSign.lean}{211}{actualLcm_integral_forces_topEdgeResidue}{top-edge residue}.}

\item\label{res:factorideal} \stag{unconditional progress}
\textbf{LCM factor-ideal shift-algebra no-go.}  For every $t\ge3$, with
$H=\lcm(1,\ldots,t)$, there is a nonzero synthetic dyadic coboundary whose
forcing letters reproduce the actual totient differences at $t-2$ prescribed
indices and lie in the ideal generated by $\ph(H)$.  Every finite integer
shift polynomial preserves the exact coboundary cancellation, the $H$-factor
ideal and all lower factor ideals, together with uniform state and letter
bounds.  A sparse-anchor variant also preserves every prescribed whole-ray
anchor.  Hence factor ideals, finite shift-algebra closure and these bounds
cannot by themselves force a contradiction.  The witness is synthetic: its
letters are not claimed to be actual totient differences, nonlinear
combinations are outside the theorem, and no cofinal family of certificates
follows.
{\footnotesize\emph{Checked:}
\mword{Erdos249257/LcmFactorIdealPulseObstruction.lean}{798}{lcm_factorIdeal_finiteRank_shiftAlgebra_not_sufficient}{finite shift-algebra no-go},
\mword{Erdos249257/LcmFactorIdealPulseObstruction.lean}{866}{lcm_factorIdeal_sparseAnchor_not_sufficient}{sparse-anchor no-go}.}

\item\label{res:periodic} \stag{unconditional progress}
\textbf{Eventually periodic weights are settled.}  For every integer base
$b\ge2$, if $\gamma\colon\N\to\Q$ is eventually periodic, nonnegative and
positive at some positive index of its periodic part, then
  $\sum_{a\ge1}\gamma(a)/(b^a-1)$ is irrational.  Theorem~A in Luca and
  Tachiya's 2017 open-access RIMS paper restates their earlier signed theorem in the exact
  purely-periodic integer case: every nonzero purely periodic integer sequence
  gives an irrational Lambert value at every integer base $|q|>1$
  \cite[Theorem~A, p.~139]{lucatachiya2017}.  Their Theorem~1 strengthens the
  nonnegative case to linear independence of each finite divisor-convolution
  ladder \cite[Theorem~1, p.~139; proof pp.~149--150]{lucatachiya2017}.
  Clearing the common denominator of a rational period and subtracting the
  finite rational prefix shows that Theorem~1 already proves the eventual
  nonnegative rational claim above; the checked theorem gives an independent
  formal proof.  By \ref{res:weight} the weight $\Aw$ is
unbounded, hence not eventually periodic, so \#249 is outside this class.
{\footnotesize\emph{Checked:}
\mword{Erdos249257/CertificateKernel.lean}{12811}{irrational_ratWeightSeries_eventuallyPeriodic}{the periodic theorem}.}

\item\label{res:primindex} \stag{proved here}
\textbf{No fixed common denominator clears the normalised weight.}  For every integer
$D\ge1$ some $n\ge1$ has $D\cdot\Aw(n)/n\notin\Z$, for $D=6$ at the prime
$n=7$, where $\Aw(7)/7=5/7$, and any $D$ clearing every coordinate up to a
horizon $N\ge4$ is divisible by an explicit two-tier primorial ($4$ at $p=2$,
$p^2$ when $p^2\le N$, else $p$).  No argument that clears the coordinates
$\Aw(n)/n$ by a single positive integer valid at all $n$ can therefore
succeed.  This obstruction is caused by the normalisation: the unnormalised
weights $\Aw(n)$ are integers.  Arguments whose common denominator grows with
the horizon, arguments using the integral weights directly, and arguments
that never clear these coordinates are untouched.
{\footnotesize\emph{Checked:}
\mword{Erdos249257/PrimitiveDeterminantLift.lean}{169}{no_fixed_integral_primitive_euler_index}{no fixed index},
\mword{Erdos249257/PrimitiveDeterminantLift.lean}{148}{twoTierPrimorial_dvd_of_integral_jet}{the primorial divisibility}.}

\item\label{res:equivalences} \stag{proved here}
\textbf{Four exact characterisations of irrationality.}  Irrationality
of $\Stot$ is equivalent to each of: non-integrality of the tail difference
$R_{N+h}-R_N$ at every $N$ and every shift $h\ge1$, equivalently a certificate
at every such pair; a certificate at every $h\ge1$ and at arbitrarily large
$N$; and the same at the scales $N=h=\lcm(1,\dots,t)$, for arbitrarily large
$t$.  A fourth equivalence,
independent of the certificate language, is a counted anti-concentration
condition on window phases.  These characterisations identify the unbounded
input exactly; they retain the full difficulty of \#249 and do not prove that
input.
{\footnotesize\emph{Checked:}
\mword{Erdos249257/LcmConeFlatness.lean}{399}{irrational_totient_series_iff_pointwise_certificates}{pointwise},
\mword{Erdos249257/LcmConeFlatness.lean}{412}{irrational_totient_series_iff_certificate_supply}{cofinal at each shift},
\mword{Erdos249257/LcmConeFlatness.lean}{426}{irrational_totient_series_iff_lcm_diagonal_certificate_supply}{$\lcm$ diagonal},
\mword{Erdos249257/PivotAntiReconstruction.lean}{1765}{dtwWindowSeparatedPairs_iff_irrational_totient_series}{window-separated pairs}.}

\item\label{res:supply} \stag{conditional reduction}
\textbf{A denominator-indexed gap-certificate hypothesis.}  Suppose that for
each denominator $q$ there is a truncation window $(N,K)$ whose totient residue
avoids a band of width $q(N{+}K{+}2)$ out of $2^K$.  This hypothesis
\emph{implies} irrationality.  It concerns a second certificate family, not
the one of \ref{res:equivalences}: the band scales with $q$, and the family is
indexed by denominators rather than by shifts.  The converse is not proved and
is not claimed.
{\footnotesize\emph{Checked:}
\mword{Erdos249257/CertificateKernel.lean}{16025}{irrational_tsum_totient_div_pow_two_of_gap_certificate_supply}{the conditional implication}.
\emph{Registry:} no row currently owns this declaration.}

\item\label{res:open} \stag{open}
\textbf{Erd\H{o}s~\#249.}  Whether $\Stot$ is irrational.  Not proved here.

\end{enumerate}
}



% === 06_replace_family_inventory_by_residue_proof; original lines 730--881 ===
% Destination: a249_family_catalogue.tex for the former family inventory; a249_front.tex thm:residueobservable for the complete pulse proof.
\section{The Lambert-series form of Problem \#249}\label{sec:family}

If $f=g*\mathbf 1$, that is $f(n)=\sum_{d\mid n}g(d)$, and
$\sum_{d\ge1}|g(d)|/(2^d-1)<\infty$, then absolute convergence justifies
interchanging the two sums and gives
\[
  \sum_{n\ge1}\frac{f(n)}{2^n}
  \;=\;\sum_{d\ge1}\frac{g(d)}{2^d-1}.
  \tag{1}\label{eq:lambert}
\]
whose right-hand side is a Lambert series with coefficients $g$.  Four
classical coefficient sequences have this shape.  The table records the
constant, the weight $g=f*\mu$, and its status.

\begin{table}[h]
\centering
\small
\begin{tabular}{@{}
  >{\raggedright\arraybackslash}p{0.07\linewidth}
  >{\raggedright\arraybackslash}p{0.19\linewidth}
  >{\raggedright\arraybackslash}p{0.21\linewidth}
  >{\raggedright\arraybackslash}p{0.43\linewidth}@{}}
\toprule
\papertablehead
$f$ & $\sum f(n)/2^n$ & weight $g=f*\mu$ & status \\
\midrule
$\tau$ & $1.6066951\ldots=\EB$ & $g\equiv1$ & The Erd\H{o}s--Borwein constant;
  irrational~\cite[theorem on p.~63]{erdos1948};
  formalised here in the Lambert form $\sum_{d\ge1}(2^d-1)^{-1}$
  (\mword{Erdos249257/CertificateKernel.lean}{8335}{irrational_erdosBorwein_series}{checked}) \\
$\omega$ & $0.5169428\ldots$ & $g=\mathbf 1_{\text{primes}}$ &
  Irrational at base $2$: Tao--Ter\"av\"ainen, Thm.~1.3, p.~4;
  proof pp.~44--56~\cite{taoteravainen2025} \\
$\Omega$ & $0.5895032\ldots$ & $g=\mathbf 1_{\text{prime powers}}$ &
  Asserted \emph{ibid.}\ by a similar argument, with the details left to the
  reader; not a proved theorem there \\
$\ph$ & $1.3676308\ldots=\Stot$ & $g=\Aw=\ph*\mu$ &
  \textbf{Open} (Erd\H{o}s \#249) \\
\bottomrule
\end{tabular}
\end{table}

\noindent The first three weights are bounded, taking only the values $0$
and~$1$ (for $\tau$ the weight is the constant $1$).  The fourth is unbounded,
by \ref{res:weight}.  That difference is the one used in
\ref{res:periodic}, and we do not claim it is the only difference: $\Aw$ also
vanishes on $n\equiv2\pmod4$, while $\mathbf 1_{\text{primes}}$ is bounded but
no more eventually periodic than $\Aw$ is, so the settled $\omega$ row is not
an instance of \ref{res:periodic} either.  It was settled by a different
method, discussed in Section~\ref{sec:open}.

Luca and Tachiya's result is stronger than the single $\tau$ row:
Example~1 makes \(1\) and every finite ladder of generalized-divisor-function
values linearly independent
~\cite[Example~1, p.~140; proof pp.~149--150]{lucatachiya2017}.
Their Example~2 applies the period-two odd-support indicator and proves joint
linear independence of its Lambert value with every finite higher
divisor-convolution ladder, for integer bases \(q\) of either sign with
\(|q|>1\) \cite[Example~2, p.~140]{lucatachiya2017}. Thus the cited theorem
strictly strengthens the isolated odd-support irrationality row.
Historically, Erd\H{o}s already singled out the totient series itself as a
difficult analogue at the end of his 1948 paper~\cite[p.~66]{erdos1948};
that remark supplies lineage, not a modern status proof.

\paragraph{A nearby 2026 totient theorem has different coordinates.}
Kaneko, Suzuki and Tachiya prove that if \(f(n)\) is a nonnegative integer
sequence of infinite support with
\(\sum_{n\le x}f(n)=O(x(\log x)^\delta)\), then, for every integer
\(t\ge2\), both
\[
  \sum_{n\ge1}\frac{f(n)}{t^{\sigma(n)}}
  \qquad\text{and}\qquad
  \sum_{n\ge1}\frac{f(n)}{t^{\ph(n)}}
\]
are irrational
\cite[Corollary~3, pp.~5--6; proof pp.~18--19]{kanekosuzukitachiya2026}.
This includes substantial families of totient-related series, but
\(\ph(n)\) occurs in the \emph{exponent}.  Problem~\#249 instead places
\(\ph(n)\) in the coefficient of \(2^{-n}\).  The theorem therefore supplies
a genuine adjacent result and a useful warning against a tempting
misidentification; it does not settle or reduce the displayed problem.

\paragraph{Exact identities and representations.}
Several neighbouring Dirichlet-convolution identities are rational and
exactly computable.  Writing $L(g):=\sum_{d\ge1}g(d)/(2^d-1)$, we have
$\sum_{d\ge1}\mu(d)/(2^d-1)=\tfrac12$
(\mword{Erdos249257/CertificateKernel.lean}{18434}{tsum_moebius_div_two_pow_sub_one_eq_half}{checked})
and $\sum_{d\ge1}\ph(d)/(2^d-1)=2$
(\mword{Erdos249257/CertificateKernel.lean}{18429}{tsum_totient_div_two_pow_sub_one_eq_two}{checked}).
These belong to two distinct one-step M\"obius-convolution chains:
\[
  L(\mathbf 1)=\EB \ \longmapsto\ L(\mu)=\tfrac12,
  \qquad
  L(\ph)=2 \ \longmapsto\ L(\ph*\mu)=\Stot.
\]
Thus $\Stot$ is one convolution step from the rational value $L(\ph)=2$, not
from $\EB$; the parallel chains show that this transformation need not
preserve rationality.  The squared-denominator representation
$\Stot=\tfrac12+\sum_{d\ge1}\mu(d)/(2^d-1)^2$
(\mword{Erdos249257/CertificateKernel.lean}{18454}{totient_series_eq_half_add_moebius_mersenne_square}{checked})
converges fast enough to recompute the decimal above.

There is also an exact probabilistic reading, and it is more than a gloss.
Let $X$ and $Y$ be independent fair-coin geometric waiting times on $\Npos$.
Then
\[
  \Stot=\tfrac12+\Pr[\gcd(X,Y)=1]
\]
(\mword{Erdos249257/CertificateKernel.lean}{18557}{totient_series_eq_half_add_visible_coprime_pairs}{checked}),
the probability being the visible-coprime-pair series
$\sum_{\gcd(m,n)=1}2^{-(m+n)}$.  That probability is related to a distribution
over the Stern--Brocot tree in a completely explicit way.  For positive
coprime $(a,b)$ put $w_{a,b}=1/(2^{a+b}-1)$.  These reduced-slope masses sum to
exactly~$1$
(\mword{Erdos249257/GcdMomentCalculus.lean}{349}{tsum_pos_coprime_inv_mersenne_eq_one}{checked}),
and at each node $(a,b)$ the cylinder mass $1/((2^a-1)(2^b-1))$ splits exactly
into its stop mass and the masses of its two children
(\mword{Erdos249257/GcdMomentCalculus.lean}{474}{cylinderMass_split}{checked}).
The coprimality probability is not that total mass; rather,
\[
  \Stot-\tfrac12
  =\sum_{\substack{a,b\ge1\\(a,b)=1}}2^{-(a+b)}
  =\sum_{\substack{a,b\ge1\\(a,b)=1}}
     \bigl(1-2^{-(a+b)}\bigr)w_{a,b}.
\]
It is therefore the expectation of the conditional primitive-pair factor
$1-2^{-(a+b)}$ under the self-similar reduced-slope probability law.  What is
missing is not structure; it is an arithmetic consequence of the structure.

The squared layer gives a different, genuinely moment-level view of the same
fair-coin model.  Since
\[
  \Pr[d\mid X\text{ and }d\mid Y]=\frac{1}{(2^d-1)^2},
\]
the checked identity
\[
  \sum_{d\ge1}\frac{\ph(d)}{(2^d-1)^2}
    =\sum_{n\ge1}(\mathsf P(n)-n)2^{-n},
  \qquad \mathsf P(n)=\sum_{e\mid n}\ph(e)(n/e),
\]
is exactly the divisor expansion of $\mathbb E[\gcd(X,Y)]$.  The companion
constant-weight identity replaces $\ph$ by $1$ and yields
$\sum_{n\ge1}(\sigma(n)-\tau(n))2^{-n}$, the corresponding divisor-count
layer.  These statements are useful because they turn the squared Lambert
denominators into observables of the same random gcd, but they do not convert
the first moment into an irrationality certificate for $\Stot$; the missing
arithmetic consequence remains the open endpoint.
{\footnotesize\emph{Checked:}
\mword{Erdos249257/GcdMomentCalculus.lean}{216}{tsum_one_div_mersenne_sq_eq_sigma_sub_tau_series}{divisor-count layer},
\mword{Erdos249257/GcdMomentCalculus.lean}{235}{tsum_totient_div_mersenne_sq_eq_gcd_moment_series}{first gcd moment}.}



% === 07_move_old_basis_history_out_of_proof; original lines 881--990 ===
% Destination: a249_p4.tex sec:mahler-defect, immediately after its kernel discussion.
\section{Proof of the dyadic-section basis theorem}\label{sec:rank}

\subsection{Prior work}

For an integer $k\ge2$, the \emph{$k$-kernel} of a sequence $c$ is
$\{n\mapsto c(k^jn+r):j\ge0,\ 0\le r<k^j\}$, and $c$ is \emph{$k$-regular}
when the module generated by its $k$-kernel is finitely generated; both
notions are due to Allouche and Shallit
(\cite[Defs.~1.1 and 2.1, author-preprint pp.~2--3]{allouche-shallit}).
Their Theorem~2.2 gives equivalent finite-kernel and matrix
characterisations, and Theorem~3.1 proves closure under convolution
~\cite[pp.~3--4 and 10--11]{allouche-shallit}; neither result controls the
finite-level rank of a sequence that is not $k$-regular.  Coons proved that
$\ph$ is not $k$-regular for any $k\ge2$ (\cite{coons}, Theorem~3.2,
pp.~348--349, in the published version; Theorem~3.3, pp.~8--9, in the
preprint, which numbers its results on a single running counter).  Bell and
Smertnig reach a further negative
conclusion, by a different argument rather than by strengthening that one:
their Theorem~1.3 shows that a $k$-Mahler series with multiplicative
coefficients has a $k$-regular coefficient sequence in an explicit closed form,
and that the totient generating series is not $k$-Mahler for any $k\ge2$
follows directly from the closed form~\cite{bell-smertnig}.

A separate line of work settles linear independence of totient values along
affine progressions, and settles it more strongly than anything proved here.
Martin~\cite{martin-phi-inequalities} assumes only that $a_1,\dots,a_m$ are
positive integers and that $a_ib_j\neq a_jb_i$ for $i\neq j$, and proves in his
Theorem~1 that for every constant $C>0$ the simultaneous ratio gaps
\[
 \frac{\ph(a_1n+b_1)}{\ph(a_2n+b_2)}>C,\quad\dots,\quad
 \frac{\ph(a_{m-1}n+b_{m-1})}{\ph(a_mn+b_m)}>C
\]
hold on a set of positive lower density; he also notes that the hypotheses are
symmetric in the forms, so every one of the $m!$ orderings occurs on such a set.

Taking $C>\bigl(\sum_{i\neq t}|c_i|\bigr)/|c_t|$ after moving a channel with
$c_t\neq0$ to the front shows at once that
\[
 \bigl\{n\mapsto\ph(a_in+b_i)\bigr\}_{i=1}^m
\]
is linearly independent over $\R$, hence over $\Q$.  His Corollary~4 carries
Theorem~1 and its corollaries to $\sigma$.  Every affine-totient independence
statement used below is therefore a consequence of Martin's theorem, and no
originality is claimed for it.

Martin does not state the exact finite-level
rank, the explicit basis, or the complete relation normal form; the present
section derives them by combining his theorem with the reductions below.
Coons and Bell--Smertnig likewise do not state those finite-level conclusions.

Coons's statement is about the union over all levels: no finite set of sections
generates the rest.  Bell and Smertnig's is about a functional equation.

One consequence of Coons's theorem should be recorded explicitly, because it
bounds what the present section may claim.  For a subgroup $M\subseteq\Z^{\N}$
of integer-valued sequences, finite generation over $\Z$ is \emph{equivalent} to
finite dimension of its $\Q$-span.  One direction is immediate.  For the other,
choose $\Q$-independent $m_1,\dots,m_d\in M$ and evaluation points
$n_1,\dots,n_d$ with $D:=\det[m_i(n_j)]\neq0$, necessarily a nonzero integer;
for $x\in M$ Cramer's rule applied to the integer system
$x(n_j)=\sum_i a_im_i(n_j)$ places every $a_i$ in $\tfrac1D\Z$, so
$M\subseteq\tfrac1D\bigl(\Z m_1+\dots+\Z m_d\bigr)$, and over the principal
ideal domain $\Z$ a submodule of a finitely generated module is finitely
generated.  (The familiar obstruction, that $\Z[\tfrac12]$ has rank one without
being finitely generated, cannot occur here: an integer-valued sequence bounds
its own denominators through its values at small indices.)

Hence \emph{infinite-dimensionality of the full totient kernel is not an
independent contribution of this note}: it is equivalent to Coons's
non-$k$-regularity and follows from it.  Only the finite-level statements are
additional.  Knowing that a span is not finitely generated says nothing
about its dimension at a given truncation, nor about which relations hold
there; what is added here is that
the dimension through level $e$ is exactly $2^e+1$, that an explicit family is
a basis, and that the two elementary reductions generate every relation.  We
found no prior source giving any of those three for $\ph$; that is the outcome
of a search, not a proof of non-existence, and no priority is claimed.

The next subsection gives an independent, constructive proof of the same
linear independence.

\subsection{Independence by a diagonal evaluation matrix}

\ref{res:rank} is proved by exhibiting a square evaluation matrix that is diagonal and invertible modulo a chosen prime.  The arithmetic construction is the whole argument.

Consider finitely many affine forms $L_i(n)=a_in+b_i$ with positive integer slopes and $a_ib_j\neq a_jb_i$ for $i\neq j$.  Write $g_i=\gcd(a_i,b_i)$ and $A_i(n)=(a_in+b_i)/g_i$.  Choose a prime $\ell$ avoiding every slope, every $g_i\ph(g_i)$, and every nonzero determinant $a_ib_j-a_jb_i$.  For a fixed row $i$ and each $j\neq i$, choose a distinct prime $q_{ij}\equiv1\pmod{\ell}$ avoiding the same finite set of divisors, and impose $L_j(n)\equiv0\pmod{q_{ij}}$ together with $A_i(n)\equiv2\pmod{\ell}$.  The Chinese remainder theorem combines the conditions; Dirichlet's theorem then supplies arbitrarily large $n$ at which $p=A_i(n)$ is prime and $p\nmid g_i$.  The diagonal value is $\ph(g_i)(p-1)\not\equiv0\pmod{\ell}$; each off-diagonal value is divisible by $\ell$.  One choice per row produces an evaluation matrix that is diagonal and invertible modulo $\ell$, so the rational determinant is nonzero
(\mword{Erdos249257/AllBaseTotientKernel.lean}{644}{linearIndependent_totientAffineForms}{checked}).

Restricting a finite relation among members of $\mathcal B$ to $n=2m+1$ makes the two zero-residue sections proportional and leaves pairwise nonproportional odd-residue forms.  The matrix argument eliminates every odd-residue coefficient; evaluation at $n=2$ separates the remaining pair, since $\ph(2)=1$ and $\ph(4)=2$.  Even residues never enter, because the second reduction has already carried them to odd residues at lower levels
(\mword{Erdos249257/TotientMahlerDefect.lean}{169}{totientKernel_even_residue_reduce}{checked}).  The same construction, with $k\nmid r$ rather than coprimality, is the independence half of Theorem~\ref{thm:kkernelrank}.  A dyadic specialisation using a power of two in place of $\ell$ is recorded as a parity-separated certificate
(\mword{Erdos249257/TotientMahlerDefect.lean}{823}{paritySeparatedMatrix_det_ne_zero}{checked},
\mword{Erdos249257/TotientMahlerDefect.lean}{882}{exists_separatedMinorCertificate_totientAffineOddFamily}{checked}).

\ref{res:basis} then needs only two further steps, and both are short.  Linear
independence is a property of finite subfamilies, so the level-$e$ statement,
applied with $e$ above the largest level occurring in a given finite subset,
already gives independence of the whole infinite family $\mathcal B$.  And
the two reductions carry every remaining section onto a rational multiple of a
member of $\mathcal B$, so the spans agree.  That the assembly is short is the
point: the level-$e$ theorem was always the whole content, and stating its
consequence as infinite-dimensionality understated it.  The sharp form is a
basis.

\begin{remark}
\ref{res:basis} and \ref{res:rank} concern $\ph$.  Their complete proofs are
independent of the arithmetic status of $\Stot$.  The single proved transport
to \#249 is \ref{res:carryrank}.
\end{remark}



% === 08_tail_equivalence_and_full_depth_proofs; original lines 990--1005 ===
% Destination: a249_p1b.tex prop:D5cons for the carry material; prop:TE-05 for the full-depth and canonical replacements.
\section{Carry-rank consequences of rationality}\label{sec:carry-rank}

Multiplying successive scaled tails by $2$ turns a hypothetical rational
value into an integral recurrence.  For coefficients $c:\N\to\N$ satisfying
$c(n)\le n$, and in particular for $c=\ph$ since $\ph(n)\le n$, the series
$\sum c(n)/2^n$ is
rational exactly when an integral scaled-tail sequence exists, that is, an
integer sequence $u$ with
$u(N+1)=2u(N)-v\,c(N+1)$ for every $N$ and $u(N)/2^N\to0$, for some integer
$v\ge1$
(\mword{Erdos249257/GenericTailOrbitRigidity.lean}{426}{binaryCoeffSeries_rational_iff_exists_temperedBinaryOrbit}{checked}).
Transporting \ref{res:rank} through that recurrence gives \ref{res:carryrank}.
The complementary finite-rank upper bound, which would make this a proof of
irrationality, is not available.



% === 09_01_transfer__section_a_farey_mediant_denominator_exclusion_lab; original lines 1005--1069 ===
% Destination: See LONG_RECORD_MAP.md, keyed by all original labels in this block.
\section{A Farey-mediant denominator exclusion}\label{sec:denominator}

\ref{res:denominator} is a Farey exclusion, and we say so before saying
anything else about it.  The argument is the classical mediant one.  The Lean
lemma commits the first $240$ binary digits of the shifted series as a single
residue over the totient carry window $(N,K)=(1,240)$; a rational $a/q$ can
equal $\Stot$ only if it falls into a resulting bad interval of width
$243/2^{240}$; that interval is bracketed by two explicit unimodular fractions
$a_1/b<c_1/d$ with $bc_1-a_1d=1$; and the mediant lemma
(\mword{Erdos249257/GapFareyBound.lean}{51}{farey_gap}{checked}) forces any
rational strictly between unimodular neighbours to have denominator at least
$b+d$.  The excluded range is therefore $q\le b+d-1$, which is the printed
constant
(\mword{Erdos249257/GapFareyBound.lean}{176}{gap_check_window_1_240_le_79639646646701375323355774875831053}{checked}),
and the transport to the series is a tail estimate
(\mword{Erdos249257/CertificateKernel.lean}{15986}{tsum_totient_div_pow_two_ne_int_div_of_totient_gap_certificate}{checked}).
\emph{The method is standard, and we claim no novelty for it.}

Two things should be read off correctly.  First, the scale.  A window of $K$
computed binary digits yields a mediant bound of order $2^{K/2}$: the same
calculation gives $2.49\times10^{17}$ at $K=120$
(\mword{Erdos249257/GapFareyBound.lean}{88}{gap_check_window_1_120_le_248672326362367909}{checked})
and $7.96\times10^{34}$ at $K=240$.  The constant is a function of how many
digits were committed, not a measure of how much the problem has moved.
Second, the sharpness.  Within its window the bound cannot be improved at all:
$b+d$ itself fails the certificate, so $b+d-1$ is exactly the last denominator
excluded.  A reader who regards an explicit Farey exclusion as routine is not
in disagreement with this note.  \ref{res:basis}, not \ref{res:denominator},
is where we would ask a sceptical reader to look.

\paragraph{The unrestricted floor.}
A certified continued-fraction prefix gives a larger exclusion by the same
kind of bookkeeping.  If $\Stot$ were rational its expansion would terminate
and its reduced denominator would be that of the last convergent, so a
certified prefix of length $n+1$ forces $q\ge q_{n+1}$.  Certified here means
that each partial quotient is emitted only when both endpoints of an exact
rational bracket $lo\le2^B\Stot\le hi$ agree on its integer part, so no
quotient is read off a float and the prefix stays genuine even if $\Stot$ is
rational.  Bracketing by the exact prefix
$lo=\sum_{n\le B}\ph(n)2^{B-n}$ with the omitted tail under
$\sum_{n>B}n2^{-n}=(B+2)2^{-B}$, and taking $B=80000$ bits, certifies $23438$
partial quotients and gives
\[
  q\ge2^{39989}\qquad\text{and}\qquad q>10^{12038},
\]
with the separation $|q_n\Stot-p_n|>0$ verified against the bracket.  At
$B=40000$ the same computation gives $10^{6018}$, so the printed figure is a
function of the bit budget.  This is an exact computation and is not
formalised; the Lean-checked exclusion of this section is the Farey instance
above.  The two are complementary rather than nested: the dyadic-family
exclusion with margin reaches denominators near $2^{200000}$ inside a
restricted family, while the continued-fraction bound excludes every
denominator below $10^{12038}$ without restriction.  The certified quotients
are Diophantine-generic, matching L\'evy's constant to $0.33$ percent with
largest partial quotient $29403$, so the series offers no anomalous convergent
to test the Mersenne machinery against.

The Farey bound transfers to the coprimality form of Section~\ref{sec:family} with
the constant halved: the visible-coprime-pair series has no representation
$a/d$ with $d\le39\,819\,823\,323\,350\,687\,661\,677\,887\,437\,915\,526$
(\mword{Erdos249257/CertificateKernel.lean}{18572}{tsum_visible_coprime_pairs_ne_int_div_of_den_le_39819823323350687661677887437915526}{checked}),
that number being $(q_0-1)/2$ for the constant $q_0$ of
\ref{res:denominator}.



% === 09_02_transfer__section_fresh_prime_losses_in_the_diagonal_probes; original lines 1069--1270 ===
% Destination: See LONG_RECORD_MAP.md, keyed by all original labels in this block.
\section{Fresh-prime losses in the diagonal probes}\label{sec:fresh-prime}

The finite diagonal and curvature probes have a useful decomposition that
keeps the effect of new prime support visible.  Fix \(H>0\) and an offset
\(s\in\N\), and write
\[
 g_H(s)=\gcd(\operatorname{rad}(H),s),
 \qquad
 M_H(k,s)=\frac{kH+s}{g_H(s)}\ph(g_H(s)).
\]
Here \(\operatorname{rad}(H)\) is the squarefree kernel of \(H\).  Since
\(g_H(s)\) divides both \(H\) and \(s\), it divides \(kH+s\).  The quantity
\(M_H(k,s)\) is the endpoint mass predicted by retaining only the prime
support already present in \(H\).  Adding the remaining prime support can only
lower the totient density, so the integer deficit
\[
 F_H(k,s)=M_H(k,s)-\ph(kH+s)
\]
is nonnegative.  This sign is the key input: the new primes are not an
uncontrolled error, even though their placement between endpoints is not
favourable by default.

Let
\[
 \Delta_H(s)=\ph(2H+s)-\ph(H+s),
 \qquad
 \Delta_H^{\rm old}(s)=M_H(2,s)-M_H(1,s).
\]
The checked endpoint calculation gives the exact identity
\[
 \Delta_H(s)=\Delta_H^{\rm old}(s)+F_H(1,s)-F_H(2,s),
\]
and the literal foreign channel is precisely the lower-endpoint deficit minus
the upper-endpoint deficit.  Thus the old-prime model is affine in the
multiplier \(k\), while the actual diagonal departs from it only through a
signed difference of two nonnegative quantities.

For the offset second difference
\(
 \delta f(s)=f(s+1)-2f(s)+f(s+2)
\)
and the five-point branch
\(Wf(s)=4\delta f(s)+\delta f(s+2)\), the same split remains exact.  Dropping
the favourable terms using \(F_H(k,s)\ge0\) gives the concrete lower bound
\[
 W\Delta_H(s)\ge W\Delta_H^{\rm old}(s)-\mathcal A_H(s),
\]
where
\[
 \mathcal A_H(s)=8F_H(1,s)+4F_H(2,s+1)+2F_H(2,s+2)
                  +F_H(2,s+3)+F_H(2,s+4).
\]
This is the hard mechanism behind the finite curvature margins: a future
cofinal argument must control the adverse budget \(\mathcal A_H(s)\).  Control
of the old-prime margin alone is insufficient.  The current formal result supplies no such cofinal
bound, so this decomposition is a finite estimate, not an irrationality proof.
{\footnotesize\emph{Checked:}
\mword{Erdos249257/FreshPrimeDeficitDecomposition.lean}{91}{endpointFreshDeficit_nonneg}{endpoint nonnegativity},
\mword{Erdos249257/FreshPrimeDeficitDecomposition.lean}{171}{diagonalHeightIncrement_eq_old_add_endpointFreshDeficits}{diagonal split},
\mword{Erdos249257/FreshPrimeDeficitDecomposition.lean}{184}{finiteForeignChannelIncrement_eq_endpointFreshDeficit_sub}{foreign-channel identity},
\mword{Erdos249257/FreshPrimeDeficitDecomposition.lean}{269}{oldFivePoint_sub_secondBranchAdverseDeficit_le_actual}{five-point loss bound}.}

\paragraph{Finite square-CRT correction suppression.}
Let \(E\) be a finite family of distinct primes \(p_i\), with anchors \(A_i\),
and let \(J<p_i\) for every \(i\in E\).  The checked square-CRT construction
produces a common base \(n\), with
\[
 n<\left(\prod_{i\in E}p_i^2\right)(\sup_{i\in E}A_i+2),
 \qquad n=A_i+p_i^2t_i
\]
for suitable \(t_i\) at each index.  Therefore every \(1\le h\le J\) obeys
\[
 \ph(n+p_i h-A_i)=(p_i-1)\ph(p_i t_i+h).
\]
The congruence has removed the \(p_i\)-divisibility correction on the whole
finite horizon, simultaneously across the family.  The formal regression
witnesses show why this is not already a separation theorem: one clean block
has both displayed coefficients zero (at base \(52\)), while another clean
block has a displayed coefficient \(-4\) (at base \(27\)).  Thus the CRT
mechanism controls a finite local error term but supplies neither a cofinal
gap nor an irrationality criterion.
{\footnotesize\emph{Checked:}
\mword{Erdos249257/SquareCRTCube.lean}{297}{exists_squareCRT_clean_totient_family}{simultaneous finite family},
\mword{Erdos249257/SquareCRTCube.lean}{327}{exists_squareCRT_clean_horizon}{finite horizon},
\mword{Erdos249257/SquareCRTCube.lean}{459}{squareCRT_clean_block_can_vanish}{vanishing block},
\mword{Erdos249257/SquareCRTCube.lean}{477}{squareCRT_clean_block_can_be_nonzero}{nonzero block}.}

\paragraph{Canonical strict-jump slack.}
At $t$, let $m_t$ be the checked canonical suffix depth and let
$d_t$ be its residue at offset $0$.  The scalar
\[
 \lambda_t=\min\!\left(d_t-2^{m_t-5},
       (2^{m_t}-2^{m_t-5})-d_t\right)
\]
is the signed distance from that residue to the nearer edge of the central
band.  The two central inequalities therefore collapse exactly to
$\lambda_t\ge0$.  The formal equivalence says that asking for this
sign at arbitrarily large strict LCM jumps is neither weaker nor stronger
than asking for the original two-sided central band there.  The same source
then deduces irrationality of $\Stot$ from the scalar condition.  The hard
step has only moved: no cofinal nonnegative-slack supply is proved, so this
section records a one-integer target for future arithmetic or computational
work, not a closure of Problem~\#249.
{\footnotesize\emph{Checked:}
\mword{Erdos249257/DiagonalFreshLossBridge.lean}{2667}{canonicalAdjacentSuffixCentral_iff_slack_nonneg}{centrality equivalence},
\mword{Erdos249257/DiagonalFreshLossBridge.lean}{2757}{canonicalAdjacentSuffixJumpCentralSupply_iff_slackSupply}{equivalence of the two conditions},
\mword{Erdos249257/DiagonalFreshLossBridge.lean}{2932}{irrational_totientSeries_of_canonicalAdjacentSuffixJumpSlackSupply}{conditional irrationality theorem}.}

\paragraph{Finite foreign-residue projection.}
For a diagonal height $H$ and cutoff $D$, retain separately the residue
channels with $d\nmid H$ and those with $d\mid H$.  The formal finite state
splits exactly as
\[
 \begin{aligned}
 \operatorname{finiteResidueDiagonal}(H,D)
   &=\operatorname{projectedForeignDefect}(H,D)\\
   &\quad+\operatorname{projectedDivisorChannels}(H,D).
 \end{aligned}
\]
Once $H\le D$, the second summand is the explicit divisor shadow.  In the
stable range $2H\le D$, the omitted foreign window has the closed bound
\[
 \left|\operatorname{foreignTailWindow}(H,D,L)\right|
 \le \operatorname{diagonalCoefficient}(H)\left(
       \frac{2}{2^D}+\frac{4}{3\,4^D}\right).
\]
The criterion is now transparent: if the limiting foreign defect is
controlled by this budget, and the finite projected state is farther from
every integer than the budget, the full target cannot be hit.  The first
comparison and a cofinal supply of separated projections are the remaining
analytic and arithmetic obligations.  No unbounded family is proved here.
{\footnotesize\emph{Checked:}
\mword{Erdos249257/ActualForeignResidueProjection.lean}{276}{abs_foreignTailWindow_le_foreignComplementBound}{finite complement bound},
\mword{Erdos249257/ActualForeignResidueProjection.lean}{308}{finiteResidueDiagonal_eq_projectedForeign_add_divisor}{exact channel split},
\mword{Erdos249257/ActualForeignResidueProjection.lean}{414}{scaleFullTarget_miss_of_projected_separation}{conditional target-miss theorem}.}

\paragraph{The actual LCM orbit: exact frontier and sign trap.}
For the actual diagonal, set $H_a=\lcm(1,\ldots,2^a)$ and
$R_a=\operatorname{totientTail}(2H_a)-\operatorname{totientTail}(H_a)$.
The formal criterion
\mword{Erdos249257/TotientActualLcmOrbitNonintegrality.lean}{37}{irrational_totientSeries_iff_actualLcmOrbitNonintegralitySupply}{is exact}
says that irrationality is equivalent to finding $R_a\notin\Z$ cofinally in
$a$.  This is the sharp endpoint reformulation: it does not ask the orbit to
stay a fixed distance from the integers, but it also supplies no such $a$.

There is nonetheless an unconditional geometric constraint on any attempted
certificate.  For $a\ge8$ and
$J+(a+6)<2\cdot2^a$, the first lookahead letter in the short LCM window is
positive and large enough to dominate the directed infinite remainder; the
whole translated difference is therefore positive.  The checked argument
uses the arithmetic positivity of the LCM-ray letters and the directed tail
bound, not an assumption about the sign of a finite residue.  In an integral
orbit this positivity is transported by the carry recurrence: the true
survivor is negative, and, once the dyadic modulus exceeds the endpoint strip,
the window residue is forced to the top edge rather than the central arc
\mword{Erdos249257/TotientActualLcmOrbitSign.lean}{211}{actualLcm_integral_forces_topEdgeResidue}{by the exact top-edge theorem}.
The sign corridor therefore removes one branch but exposes the remaining
arithmetic task.  A proof must exclude this persistent top boundary (or find
cofinal non-integrality); a sign-only modular kill is not available.

\paragraph{A quantitative separation criterion.}
The adjacent arithmetic theorem makes the finite evidence at the actual
LCM heights explicit rather than treating it as a numerical table:
\mword{Erdos249257/TotientActualLcmShortKill.lean}{18}{lcmDiagonalArithmeticKill_two_pow_four}{the \(t=2^4\), length-\(23\) kill}
and
\mword{Erdos249257/TotientActualLcmShortKill.lean}{27}{lcmDiagonalArithmeticKill_two_pow_six}{the \(t=2^6\), length-\(93\) kill}
give
\mword{Erdos249257/TotientActualLcmShortKill.lean}{35}{actualLcmTailOrbit_four_and_six_notMem_int}{actual-orbit non-integrality at both heights}.
The same file records the latter witness as a short arithmetic result
through exponent six
\mword{Erdos249257/TotientActualLcmShortKill.lean}{54}{powerTwoActualLcmShortArithmeticKillSupply_through_six}{but explicitly stops at that finite prefix}; it does not supply the missing cofinal quantifier.

A stronger unbounded approach asks for a quantitative anti-concentration
bound.  A sign alone gives too little information.  At odd rank \(2q+1\), the exact raw
LCM block approximates the actual orbit with the elementary error bound
\mword{Erdos249257/TotientActualLcmOrbitSeparation.lean}{141}{abs_actualLcmTailOrbit_sub_rawApprox_lt}{recorded here};
the raw block itself is an integer half-word divided by \(4^q\)
\mword{Erdos249257/TotientActualLcmOrbitSeparation.lean}{179}{actualLcmRawApprox_eq_half_div_fourPow}{by the normalization identity}.
Separation of that raw value from every integer by \(1/32\) forces the
half-word band
\mword{Erdos249257/TotientActualLcmOrbitSeparation.lean}{208}{halfWordBandAt_of_rawApprox_integerSeparation}{through the exact threshold lemma}.
The formal cofinal hypothesis strengthens this to \(1/32\) plus the orbit
error. This condition is
\mword{Erdos249257/TotientActualLcmOrbitSeparation.lean}{254}{PowerTwoActualLcmOrbitSeparationSupply}{the cofinal separation hypothesis},
and it yields irrationality by
\mword{Erdos249257/TotientActualLcmOrbitSeparation.lean}{305}{irrational_totientSeries_of_actualLcmOrbitSeparationSupply}{the conditional endpoint theorem}.
Thus the hard step is an actual cofinal anti-concentration estimate; finite
separation, the sign corridor, and the exact endpoint reduction do not imply
it.

There is also a precise no-go for a tempting substitute.  In the positive
short LCM window, a terminal dyadic staircase would force the complete
discrepancy word to vanish modulo its terminal power of two, but the last
positive arithmetic letter is strictly smaller than that modulus, so
\mword{Erdos249257/TotientActualLcmTopEdgeStaircase.lean}{251}{not_actualLcmTerminalDyadicStaircase_of_room}{the staircase is impossible under the stated room hypothesis}.
The companion
\mword{Erdos249257/TotientActualLcmTopEdgeStaircase.lean}{297}{windowDiscrepancy_emod_two_pow_eq_terminal}{terminal-remainder identity}
shows what survives: only the final letters control that residue.  This
prunes total dyadic annihilation, not every top-edge pattern; the remaining
endpoint problem is still to prove a nonzero cofinal residue gap.



% === 09_03_transfer__section_obstructions_for_fixed_coordinate_methods; original lines 1270--1312 ===
% Destination: See LONG_RECORD_MAP.md, keyed by all original labels in this block.
\section{Obstructions for fixed-coordinate methods}\label{sec:nogo}

\ref{res:primindex} has a one-line proof.  Take a prime
$p>\max(D,2)$; then $\Aw(p)=p-2$ by \ref{res:weight}, and $p\nmid D(p-2)$, so
$D\cdot\Aw(p)/p\notin\Z$.

Two further checked results delimit narrower abstract hypothesis classes.

A fixed-precision transport theorem
(\mword{Erdos249257/TropicalCurvatureCarry.lean}{137}{fixedPrecisionTropicalNoGo}{checked})
gives an unconditional reusable negative for a whole local escape strategy:
at every fixed positive precision, every finite odd-unit valuation word admits
a prefix-locked centred completion.  The hard step is one-step modular
recentering: the compatible carry is chosen by reducing the shifted state into
an interval-sized power-of-two window, with an unrestricted quotient.  Induction
then centres every letter of the finite word, whatever the incoming carry.
Consequently, no contradiction can be extracted from a finite local
valuation-unit signature read at a precision fixed in advance; this retires
that fixed-positive-precision local strategy for good.  The theorem is
synthetic and finite-word only: its valuation/unit symbols and unrestricted
high quotients do not encode the actual totient differences, so it mentions
neither $\ph$ nor $\Stot$.  It does not provide a growing-precision or global
correlation bridge, nonlinear arithmetic, or a genuinely fresh divisor channel,
and it does not settle the \#249 endpoint or irrationality.

The factor-ideal no-go of \ref{res:factorideal}
(\mword{Erdos249257/LcmFactorIdealPulseObstruction.lean}{798}{lcm_factorIdeal_finiteRank_shiftAlgebra_not_sufficient}{checked},
with a sparse-anchor companion at
\mword{Erdos249257/LcmFactorIdealPulseObstruction.lean}{866}{lcm_factorIdeal_sparseAnchor_not_sufficient}{checked})
exhibits, for each $t\ge3$ with $H=\lcm(1,\dots,t)$, a nonzero integer carry
whose forcing letters lie in the ideal generated by $\ph(H)$, which reproduces
the true totient differences at $t-2$ prescribed indices, and which every
finite integer shift polynomial carries to a pair with the same coboundary
form, the same ideal memberships and an $\ell^1$-weight bound.  It does close
that hypothesis class.  Two limits.  The witness is a spike, whose state is
$-\ph(H)$ at $t-2$ points and zero elsewhere, so it stays uniformly bounded
while the strip bounds it respects grow, and the construction is cheap for
that reason.  And the shifted pairs are not asserted to be nonzero, so an
argument that additionally demanded non-vanishing after the shift is not
excluded.  The theorem contains no whole-ray anchor condition, diagonal bound,
or strict-survivor condition.



% === 09_04_transfer__section_conditional_criteria_label_sec_conditiona; original lines 1312--1471 ===
% Destination: See LONG_RECORD_MAP.md, keyed by all original labels in this block.
\section{Conditional criteria}\label{sec:conditional-criteria}

The criteria below do not enter the proofs of the unconditional results.
Each states its unproved hypotheses explicitly.

\paragraph{A prime-orbit sufficient condition.}
The hypothesis
\lword{Erdos249/TotientStrictPrimeEscape.lean}{18}{DTWNaturalPrimeTailOrbitStrictGap}{asks}
for cofinally many primes $p$ at which the first tail-orbit exponential has
real part strictly below $9/10$.  Granting it,
\lword{Erdos249/TotientStrictPrimeEscape.lean}{27}{naturalPivotPointEscape_of_naturalPrimeTailOrbitStrictGap}{a positive adaptive truncation budget}
follows, and then
\lword{Erdos249/TotientStrictPrimeEscape.lean}{66}{irrational_totient_series_of_naturalPrimeTailOrbitStrictGap}{$\Stot$ is irrational}.
The $9/10$ hypothesis is weaker than the earlier $4/5$ hypothesis, but no
instance of this cofinal hypothesis is proved.  Those three declarations are
the entire formal content of the module; everything else in this paragraph is
ordinary mathematics, and is labelled as such.

The remainder of this branch is elementary phase dynamics.  If
\(\zeta_{h,M}=\operatorname{tailOrbitFirstExp}(h,M)\), then
\[
  \zeta_{h,M+1}=\zeta_{h,M}^{\,2}
\]
because the integer increment in the tail-difference recurrence disappears
after exponentiation.  Iteration gives the closed form
\[
  \zeta_{h,M+k}=\zeta_{h,M}^{\,2^k},
  \qquad \zeta_{h,N}=\zeta_{h,0}^{\,2^N},
\]
and the prime-gap predicate is exactly the same $9/10$ test on these
initial-phase iterates along the shifted primes.  None of these three
observations is formalised.

The same dynamics gives an obstruction: if an iterate is a dyadic root of
unity, repeated squaring reaches the absorbing phase $1$, so the cofinal strict
gap is impossible.  A conditional bridge from phase information to the
prime-index hypothesis is plausible in the same coordinates, namely that
cofinally many nonpositive phases at shifted prime indices would imply the
strict prime-orbit gap; it is recorded here as an expectation and is neither
proved in this paper nor formalised, and in any case it would not supply those
prime-aligned visits.
These are classification and obstruction statements for one conditional
family, not a classification of the actual initial phase.
This identifies the missing input as phase-density or prime-index supply for
that squaring orbit, not another finite certificate theorem.  The identity
and its root obstruction supply neither anti-concentration, density, nor a
cofinal prime gap, so no irrationality conclusion follows.
Section~\ref{sec:frontier} states the sharper form
in which this hypothesis is now available.

\paragraph{Finite Euler-factor identities.}
The local coefficients are
\[
  \eta(0)=1,\qquad \eta(1)=-2,\qquad \eta(2)=1,\qquad
  \eta(e)=0\quad(e\ge3),
\]
the coefficients of $(1-X)^2$.  For a prime $p$ and $e\ge0$, write
$\sigma_p(e)=1+p+\cdots+p^e$.  The two local Euler factors are
\[
  1-\frac2p+\frac1{p^2}=\left(1-\frac1p\right)^2,
  \qquad
  1-\frac2{p^2}+\frac1{p^4}
      =\left(1-\frac1{p^2}\right)^2
\]
at
\lword{Erdos249/FiniteEulerSieve.lean}{28}{muSqEulerFactor_one}{$s=1$}
and
\lword{Erdos249/FiniteEulerSieve.lean}{36}{muSqEulerFactor_two}{$s=2$}.
On the prime-power divisor-sum row,
\[
  \sigma_p(1)-2\sigma_p(0)=p-1
\]
(\lword{Erdos249/FiniteEulerSieve.lean}{44}{sigmaPrimePow_firstDiff}{first difference}),
and for every $e\ge0$,
\[
  \sigma_p(e+2)-2\sigma_p(e+1)+\sigma_p(e)
    =p^{e+1}(p-1)
\]
(\lword{Erdos249/FiniteEulerSieve.lean}{51}{sigmaPrimePow_secondDiff}{second difference}).
For example, at $p=3$ the values
$\sigma_3(0),\ldots,\sigma_3(3)=1,4,13,40$ give
$4-2=2$ and $40-2\cdot13+4=18=3^2(3-1)$.  Thus convolution with
$\mu\ast\mu$ converts the prime-power divisor-sum row into the corresponding
totient row at every finite stage.  These are finite algebraic identities; no
transcendence conclusion is attached to them.

\paragraph{Mixed differences of cyclotomic layers.}
For a function $F:\{0,1\}^2\to\Z$, define
\[
  \Delta F=F(1,1)-F(1,0)-F(0,1)+F(0,0)
\]
(\lword{Erdos249/PrimeRayCyclotomicCurvature.lean}{22}{checkerboard}{mixed difference}).
If $F(i,j)=u(i)+v(j)$, then $\Delta F=0$
(\lword{Erdos249/PrimeRayCyclotomicCurvature.lean}{26}{checkerboard_separable}{separable case}).
Conversely, among integer linear combinations of the four values of $F$ whose
coefficients sum to zero along each row and each column, the pattern
$(1,-1,-1,1)$ is forced up to scale
(\lword{Erdos249/PrimeRayCyclotomicCurvature.lean}{32}{checkerboard_unique}{fixed-stencil uniqueness}).
For example, the table
\[
\begin{array}{c|cc}
 &0&1\\ \hline
0&9&13\\
1&12&16
\end{array}
\]
has mixed difference $16-12-13+9=0$; increasing only the lower-right entry by
$5$ changes the mixed difference to $5$.  The uniqueness statement is
internal to this fixed $2\times2$ stencil.  It does not assert invariance under
larger stencils, nonlinear functionals, or another representation.

The formal development also records
\lword{Erdos249/PrimeRayCyclotomicCurvature.lean}{45}{fourPoint_layer_identity}{a four-point divisor-layer identity}
and
\lword{Erdos249/PrimeRayCyclotomicCurvature.lean}{62}{residueClass_step_has_centred_completion}{a centred representative for a residue class modulo an even modulus}.
The two unproved hypotheses and their target conclusion are as follows.  For
$C:\N\to\N$ and $m,d\in\N$, put
\[
\begin{aligned}
\operatorname{Layer}(C,m)\quad\Longleftrightarrow\quad&
\exists Q_0\ \forall q\ge Q_0,\quad q\ {\rm prime}\Longrightarrow\\[-2pt]
&\hspace{5em}1<C(mq)\ \hbox{ and }\ \gcd(C(mq),mq)=1.
\end{aligned}
\]
\[
\begin{aligned}
\operatorname{BoundedOrder}(C,m,d)\quad\Longleftrightarrow\quad&
\forall q,p,\quad q,p\ {\rm prime},\ p\mid C(mq)\Longrightarrow\\[-2pt]
&\hspace{3em}\exists k,\quad 1\le k\le d\ \hbox{ and }\ mq\mid p^k-1.
\end{aligned}
\]
\[
\begin{aligned}
\operatorname{FinitePrimeEscape}(C,m)\quad\Longleftrightarrow\quad&
\forall\text{ finite sets }S\text{ of primes},\ \exists Q_0\
\ \forall q\ge Q_0,\\[-2pt]
&\hspace{3em}q\ {\rm prime}\Longrightarrow
\forall p\in S,\quad p\nmid C(mq).
\end{aligned}
\]
These are the linked predicates
\lword{Erdos249/PrimeRayCyclotomicCurvature.lean}{95}{PrimeRayLayerSupply}{Layer},
\lword{Erdos249/PrimeRayCyclotomicCurvature.lean}{102}{BoundedDegreeOrderConsumer}{BoundedOrder},
and
\lword{Erdos249/PrimeRayCyclotomicCurvature.lean}{151}{FinitePrimeSupportEscape}{FinitePrimeEscape}.
The bounded-order condition does not assert an exact multiplicative order; it
asserts only the displayed divisibility for some $k\le d$.

The implication from bounded order to finite-prime escape is a size argument
in the paper.  If $m=0$, the bounded-order condition forces $C(0)=1$.  If
$m>0$, the assertion is immediate for $S=\varnothing$; otherwise choose
$Q_0$ so that $mQ_0>\max_{p\in S}(p^d-1)$.  For every prime
$q\ge Q_0$, a divisor $p\in S$ of $C(mq)$ would give
$mq\mid p^k-1$ for some $k\le d$, hence
$mq\le p^k-1\le p^d-1$, a contradiction.  The layer hypothesis separately
ensures that the layers are nontrivial and coprime to their indices.  Neither
hypothesis is proved here; polynomial-resultant realisability and
Archimedean growth are also not established.



% === 09_05_transfer__section_open_problems_label_sec_frontier_; original lines 1471--1870 ===
% Destination: See LONG_RECORD_MAP.md, keyed by all original labels in this block.
\section{Open problems}\label{sec:frontier}

Write $H_t=\lcm(1,\dots,t)$
(\mword{Erdos249257/CarrySurvivorExtinction.lean}{515}{periodLcm}{definition}),
so that $H_1=1$, $H_2=2$, $H_3=6$, $H_4=12$, $H_5=H_6=60$ and $H_7=420$; and
recall from Section~\ref{sec:results} the tail
$R_N=\sum_{m\ge1}\ph(N+m)/2^{m}$, which differs from $2^N\Stot$ by an integer
(\mword{Erdos249257/TotientTailPeriodKiller.lean}{57}{totientTail}{definition}).

\subsection{Exact equivalent formulations}

\begin{equivform}[lcm-diagonal escape]\label{prob:diagonal}
For every $t_0$ there is a $t\ge t_0$ with $R_{2H_t}-R_{H_t}\notin\Z$.
\end{equivform}

\noindent By \ref{res:equivalences}, formulation~\ref{prob:diagonal} is
equivalent to irrationality of $\Stot$.  It identifies an exact decidable
witness family, but is not claimed to reduce the difficulty of Erd\H{o}s
\#249.  A finite list of successful certificates establishes the predicate
only on its tested scales.  Conversely, one pair $h>0$, $N$ with
$R_{N+h}-R_N\in\Z$ already makes $\Stot$ rational
(\mword{Erdos249257/LcmConeFlatness.lean}{357}{not_irrational_of_tail_diff_mem_int}{checked}).

The certificate is decidable at each $t$, and is kernel-checked at the $28$
listed scales between $t=1$ and $t=64$
(\mword{Erdos249257/DiagonalPincerCertificatesT64.lean}{1933}{diagonalPincerCertificateScalesThroughT64}{list},
\mword{Erdos249257/DiagonalPincerCertificatesT64.lean}{1967}{certifiedKill_diagonal_all_imported_through_t64}{verification}).
The $28$ certified scales through $t=64$ are the whole of the checked
coverage; no scale beyond $64$ is verified, and the finite list discharges no
part of the cofinal quantifier.

The cyclotomic form exposes more arithmetic structure.  Write
$C(n)=|\Phi_n(2)|$.

\begin{equivform}[cyclotomic-anchored phase escape]\label{prob:cyclotomic}
For every $h\ge1$ and every $N_0$, there exist primes $q,p$ and $L\in\N$ such
that
\[
 \gcd(p,hq)=1,\qquad p\mid C(hq),\qquad hq\mid p-1,
 \qquad p-1\ge N_0,
\]
and $(hq,p-1,L)$ is a finite tail-difference certificate.
\end{equivform}

\noindent This is
\lword{Erdos249/CyclotomicAnchoredKill.lean}{3231}{CyclotomicAnchoredKillSupply}{CyclotomicAnchoredKillSupply}.
For the binary layers, clean anchors satisfying every condition before the
certificate are already supplied, and
\lword{Erdos249/CyclotomicAnchoredKill.lean}{3327}{binaryCyclotomicAnchoredKillSupply_iff_irrational}{the checked equivalence}
identifies formulation~\ref{prob:cyclotomic} with irrationality of $\Stot$.
The unknown is phase escape of the totient discrepancy modulo $2^L$, not
large prime divisors or multiplicative order.  The checked $p=331$ instances
are finite models of this exact predicate; unbounded support alone does not
imply it.

The
\lword{Erdos249/CyclotomicAnchoredKill.lean}{1706}{exists_unbounded_binaryCyclotomicSupport_with_periodLock_of_not_irrational}{checked conditional support/period no-go}
makes that limitation exact.  If $\Stot$ were rational, its eventual integral
tail-period law would supply some $h>0$, while
\lword{Erdos249/PrimeRayCyclotomicCurvature.lean}{158}{UnboundedPrimeDivisorSupply}{unbounded prime divisors}
would persist on the same $h$-ray for the actual binary cyclotomic layers, by
\lword{Erdos249/CyclotomicAnchoredKill.lean}{1693}{binaryCyclotomicLayer_unboundedPrimeDivisorSupply}{the checked binary-layer theorem}.
Consequently, unbounded support cannot be the missing contradiction unless an
additional carry- or phase-escape theorem relates that support to the period
law.  This is a necessary-state theorem under hypothetical rationality: it
neither constructs a rational example nor proves that the rationality
hypothesis is consistent, and Erd\H{o}s \#249 remains open.

\subsection{Exact full-depth equivalences}

The earlier condition
\lword{Erdos249/PeriodMultipleEscape.lean}{441}{ApFullDepthEscape}{ApFullDepthEscape}
asks, for every $d\ge1$ and $N$, for a $t\ge1$ such that $(td,N,td)$ is a
finite tail-difference certificate.  It was originally recorded only as a
sufficient condition because the certificate depth is forced to equal its
period.  The full-depth theorem removes that asymmetry completely.
For fixed $d,N$, one certificate $(d,N,L)$ at any depth yields a threshold
$T>0$ beyond which every adjacent pair of multipliers contains an $m$ with a
full-depth certificate $(md,N,md)$
(\lword{Erdos249/FullDepthRayAmplifier.lean}{364}{eventually_twoSyndetic_fullDepthKillMultipliers_of_seed}{checked}).
Together with finite-certificate faithfulness, this gives the pointwise
equivalence
\[
 \exists t>0:\ (td,N,td)\text{ is certified}
 \quad\Longleftrightarrow\quad R_{N+d}-R_N\notin\Z
\]
(\lword{Erdos249/FullDepthRayAmplifier.lean}{388}{exists_fullDepthKill_on_ray_iff_shift_notMem_int}{checked}).
Therefore $\operatorname{ApFullDepthEscape}$ is equivalent to irrationality
(\lword{Erdos249/FullDepthRayAmplifier.lean}{406}{apFullDepthEscape_iff_irrational}{checked}).
The cofinal full-depth supply is likewise equivalent to the arbitrary-depth
period-multiple supply and hence to irrationality
(\lword{Erdos249/FullDepthRayAmplifier.lean}{421}{cofinalFullDepthKillSupply_iff_periodMultipleKillSupply}{depth deletion};
\lword{Erdos249/FullDepthRayAmplifier.lean}{438}{cofinalFullDepthKillSupply_iff_irrational}{endpoint equivalence}).
What remains open is the initial nonintegrality: the theorem converts one
non-integral shift, equivalently one seed, into a dense ray of full-depth
certificates but does not furnish that seed.

\begin{problem}[first-harmonic anti-concentration on supplier fibres]
\label{prob:firstharmonic}
For each $h\ge1$, do there exist $s\ge1$ and $0<\eta<1$ such that, for every
$X_0$, there are $X,L$ with
\[
 \max(X_0,1)\le X,\qquad h\le L-s,\qquad
 16(2X+h+L+2)\le2^L,
\]
for which the four bounds below hold?
\end{problem}

\noindent This four-term decomposition is one way to prove a fixed-full-block
first-harmonic gap.  The underlying question is whether the first additive
character of the totient-window discrepancy has a constant saving.  Failure
of this particular decomposition would not decide that question, still less
Erd\H{o}s Problem~\#249.

To make the question self-contained, put
\[
 \omega_N=\exp\!\left(2\pi i\,
   \frac{D_{h,N,L}\bmod 2^L}{2^L}\right),\qquad X\le N<2X.
\]
At the pivot argument $N+L-s+1$, call $N$ a supplier when this integer factors
as $mp$, where $p$ is its largest prime factor,
$0<m\le\sqrt X/2$, and $p>2\sqrt X$.  Suppliers are good when
$\eta m\le\ph(m)$ and bad otherwise; non-suppliers form the remaining set.
Writing $t=L-s+1$, the canonical fibre with cofactor $m$ is exactly the image
of the primes in
\[
 X+t\le mp<2X+t
\]
under the injective map $p\mapsto mp-t$
(\mword{Erdos249257/FirstHarmonicPivot.lean}{299}{mem_pivotFiber_iff_exists_supplierPrime}{membership},
\mword{Erdos249257/FirstHarmonicPivot.lean}{367}{pivotBaseOfPrime_injective_on_supplierPrimes}{injectivity},
\mword{Erdos249257/FirstHarmonicPivot.lean}{384}{image_pivotSupplierPrimes_eq_pivotFiber}{image equality}).
This exact bijection does not imply that the supplier prime is isolated from
every other argument in the window.  At $X=16$, $L=20$, $s=1$, $m=2$, and
$N=18$, the pivot is $38=2\cdot19$, while the same prime $19$ divides the
distinct argument $N+1=19$
(\mword{Erdos249257/FirstHarmonicPivot.lean}{394}{supplierPrime_not_globally_isolated_counterexample}{checked counterexample}).

For a supplier write $z_N$ for its explicit totient pivot phase,
$w_N=\omega_N/z_N$, and $\bar z_m$ for the mean of $z_N$ on the fibre with
cofactor $m$.  Define
\[
\begin{aligned}
 \mathcal C&=\sum_{N\ {\mathrm{good}}}w_N(z_N-\bar z_m),&
 \mathcal M&=\sum_{N\ {\mathrm{good}}}w_N\bar z_m,\\
 \mathcal B&=\sum_{N\ {\mathrm{bad}}}\omega_N,&
 \mathcal U&=\sum_{N\ {\mathrm{nonsupplier}}}\omega_N.
\end{aligned}
\]
The exact identity is
$\sum_{X\le N<2X}\omega_N=\mathcal C+\mathcal M+\mathcal B+\mathcal U$, and
Problem~\ref{prob:firstharmonic} asks precisely for
\[
 \operatorname{Re}\mathcal C\le\frac{14}{25}X,\qquad
 |\mathcal M|\le\frac1{100}X,\qquad
 |\mathcal B|\le\frac1{100}X,\qquad
 |\mathcal U|\le\frac8{25}X.
\]
The decomposition is
\mword{Erdos249257/FirstHarmonicPivot.lean}{514}{windowFirstExp_sum_eq_pivot_decomposition}{checked},
and the four budgets imply irrationality by
\mword{Erdos249257/FirstHarmonicPivot.lean}{577}{irrational_totient_series_of_pivotResidualDecorrelation}{the conditional irrationality theorem}.

More precisely, these four inequalities are the source predicate
\mword{Erdos249257/FirstHarmonicPivot.lean}{543}{PivotBudgetAt}{the budget condition},
and its checked consequence
\mword{Erdos249257/FirstHarmonicPivot.lean}{549}{first_harmonic_gap_of_pivotBudgetAt}{the first-harmonic gap}
is the explicit estimate
\[
 \sum_{X\le N<2X}\cos\!\left(2\pi
   \frac{D_{h,N,L}\bmod 2^L}{2^L}\right)\le\frac9{10}X.
\]
The hard missing step is therefore cofinal centred-correlation
and error bounds, not the formal addition of the four terms.  In particular,
the exact supplier-fibre bijection above does not provide global prime
isolation; the $38=2\cdot19$ witness is the natural friction that the
correlation estimate must overcome.  We claim neither the budget predicate nor
the resulting gap without those cofinal bounds.
An obstruction showing that available first- and second-moment information
cannot force this modulo-$2^L$ small-ball estimate would also decisively close
this harmonic argument.

\subsection{Structural and Diophantine routes}

\begin{problem}[totient carry-rank compression]\label{prob:carryrank}
Suppose $v\ge1$ and $u:\N\to\Z$ satisfy
\[
 u(N+1)=2u(N)-v\ph(N+1),\qquad \frac{u(N)}{2^N}\longrightarrow0.
\]
Must the span of the level-$e$ dyadic sections of $u$ have dimension less
than $2^e-1$ for some $e$?
\end{problem}

\noindent Rationality of $\Stot$ would produce exactly such an integral
tempered orbit, while the checked transport theorem forces its rank to be at
least $2^e-1$ at every level.  Thus a positive answer proves irrationality.
A negative answer should ideally construct a control recurrence with the same
integrality and decay architecture and near-maximal dyadic rank; infinite
dimension of the totient $2$-kernel alone is not enough.

For a separate denominator-compression argument, define the M\"obius--Mersenne
ladder
\[
 \Theta_r=\sum_{d\ge1}\frac{\mu(d)}{(2^d-1)^r},
 \qquad \Theta_2=\Stot-\frac12,
\]
whose second rung is Fan's identity~\cite{fan2026totient}.
Result~\ref{res:rankonefloor} excludes the positive rank-one family by an
exact extremal calculation, not a coarse interval screen.  The unique minimum
is $(e,Y)=(1,5)$, and every admissible quotient, hence every nonempty finite
positive normalized average of them, satisfies
\[
 Q(e,Y)-\Theta_2>\frac{21}{320}>\frac1{16}.
\]
The stronger unit-fraction floor $1/15$ fails already at the minimizer.  Thus
adding atoms or positive rank-one blocks cannot create a small linear form;
any remaining argument must use signed cancellation, off-diagonal coupling, or
genuinely higher rank.

\begin{problem}[genuinely coupled denominator compression]\label{prob:hankel}
Construct integers $P_n$ and $Q_n>0$ from a genuinely off-diagonal,
vector-valued, or signed-cancellation Hankel--Schur--Pad\'e kernel such that
\[
 0<Q_n\left|\Theta_2-\frac{P_n}{Q_n}\right|\longrightarrow0.
\]
Alternatively, prove a structural no-go theorem for a clearly defined larger
class of bounded-width coupled kernels.
\end{problem}

\noindent Ordinary convergence of rational approximants is insufficient; the
scaled error is the Diophantine quantity.  The words ``genuinely coupled''
exclude the positive rank-one and positive direct-sum mechanisms already
separated from $\Theta_2$ by the uniform gap.

\subsection{Independent extension}

This subsection previously posed the odd-prime kernel dimension as an open
problem.  It is now a theorem, for every integer base rather than for odd
primes only, and the rank, basis, and relation dimension are Lean-checked
unconditionally.  Martin's positive-density theorem remains prior art for the
independence conclusion; it is not a Lean dependency.

\begin{theorem}[exact rank of every totient $k$-kernel truncation]
\label{thm:kkernelrank}
Let $k\ge2$ and $e\ge1$ be integers, write $F^{(k)}_{j,r}(n)=\ph(k^jn+r)$, and put
\[
 V_{k,e}=\operatorname{span}_{\Q}
 \{\,F^{(k)}_{j,r}:0\le j\le e,\ 0\le r<k^j\,\}.
\]
Then $\dim_{\Q}V_{k,e}=k^e+1$, and
\[
 \mathcal B_{k,e}
 =\{F^{(k)}_{0,0},F^{(k)}_{1,0}\}
 \cup\{\,F^{(k)}_{j,r}:1\le j\le e,\ 1\le r<k^j,\ k\nmid r\,\}
\]
is a basis.  Every omitted section reduces to a basis element by an explicit
scalar: $F^{(k)}_{j,0}=k^{j-1}F^{(k)}_{1,0}$ for $j\ge1$, and if $r=k^tu$ with
$t=\max\{s:k^s\mid r\}\ge1$ and $k\nmid u$, then
\[
 F^{(k)}_{j,r}=C_k(t,u)\,F^{(k)}_{j-t,u},
 \qquad
 C_k(t,u)=k^t\prod_{\substack{p\mid k\\ p\nmid u}}\Bigl(1-\tfrac1p\Bigr).
\]
\end{theorem}

The checked composite-base statement is deliberately cross-multiplied by
$\ph(\gcd(k,u))$; the displayed scalar is the equivalent paper-level
normalisation after exact division.  Lean checks the zero-residue relation,
the one-step reduction, spanning of the complete truncation, affine-form
independence, exact rank $k^e+1$, an explicit basis, and the relation-module
dimension in \texttt{Erdos249257/AllBaseTotientKernel.lean}:
\mword{Erdos249257/AllBaseTotientKernel.lean}{767}{allBaseTotientKernel_zero_residue}{zero residue},
\mword{Erdos249257/AllBaseTotientKernel.lean}{783}{allBaseTotientKernel_step}{one-step reduction},
\mword{Erdos249257/AllBaseTotientKernel.lean}{957}{span_allBaseThroughLevelFamily_eq}{span},
\mword{Erdos249257/AllBaseTotientKernel.lean}{644}{linearIndependent_totientAffineForms}{affine independence},
\mword{Erdos249257/AllBaseTotientKernel.lean}{1231}{finrank_allBaseTotientKernel_eq}{rank},
\mword{Erdos249257/AllBaseTotientKernel.lean}{1248}{allBaseTotientKernelBasis}{basis},
\mword{Erdos249257/AllBaseTotientKernel.lean}{1344}{finrank_allBaseRelationModule_eq}{relation dimension}.
The older conditional module \texttt{TotientKernelConditional.lean} is not
the present authority.  Martin's Theorem~1 still subsumes the independence
conclusion as prior art and is not formalised; it is not a remaining Lean gap.

The selection condition is $k\nmid r$, not $\gcd(k,r)=1$; for composite $k$ a
basis residue may share prime factors with $k$.  Independence is obtained by
restricting to $n=km+1$, where the surviving channels become the affine forms
$L_0(m)=km+1$ and $L_{j,r}(m)=k^{j+1}m+(k^j+r)$ with $k\nmid r$.  These have
positive slopes and satisfy $a_ib_j\neq a_jb_i$.  Martin's
Theorem~1~\cite{martin-phi-inequalities} subsumes the independence conclusion
as prior art.  Lean discharges it by the finite-determinant row
\texttt{linearIndependent\_totientAffineForms}, not by importing Martin's density theorem.  The two zero-residue channels
are proportional on that progression and are separated afterwards by evaluating
at $n=k$, using $\ph(k^2)=k\ph(k)$ and $\ph(k)<k$.

At $k=2$ this recovers the dyadic rank theorem of Section~\ref{sec:rank}; at
$k=\ell$ an odd prime it gives the case first posed here.  Neither Coons's
non-$k$-regularity theorem nor the Bell--Smertnig Mahler classification supplies
a finite-level rank, a basis, or the relations, so
Theorem~\ref{thm:kkernelrank} is not a consequence of either.  A targeted
literature search located no source stating the exact rank, the basis, or the
normal form; that is a search result and not a novelty claim, and the statement
remains subject to specialist review.

\begin{remark}[excluded generic mechanisms]
Unbounded prime support, exact multiplicative order, natural-boundary
information, ordinary Hankel nonvanishing, and positive rank-one averaging do
not supply the missing binary phase or denominator compression.  The finite
Euler-factor and $2\times2$ mixed-difference identities of
Section~\ref{sec:conditional-criteria} are tools rather than open endpoints; no
growing-level Diophantine conclusion is claimed for them here.
\end{remark}

A proof of formulation~\ref{prob:diagonal} gives irrationality of $\Stot$ by
\ref{res:equivalences}.  Its negation settles \#249 in the opposite direction:
it supplies a positive-shift integral tail difference, which forces
rationality.  By contrast, failure of the sufficient bounds would rule out
only those particular arguments.

\paragraph{The independent denominator exclusion.}
The constant $7.96\times10^{34}$ of \ref{res:denominator} is the classical
Farey mediant bound applied to a window with a free parameter $K$; it has
order $2^{K/2}$, and raising it needs only more committed binary digits, which
is a computation rather than an idea.  The unrestricted floor
$q\ge2^{39989}$ and $q>10^{12038}$ of Section~\ref{sec:denominator} is larger and is
the same kind of statement, bought with more bits rather than more structure.
Neither is one of the open items of Section~\ref{sec:open}.

\subsection{Logical status and analytic input}\label{sec:open}

\begin{enumerate}
\item Irrationality of $\Stot$ (\ref{res:open}).  No proof is claimed.
\item The lcm-diagonal and clean cyclotomic-anchor statements are exact
  equivalent formulations of \ref{res:equivalences}: they convert
  irrationality into the existence of an unbounded family of finite
  tail-difference certificates.  Finite instances are checked, every
  certified scale $t\le64$, the $28$-scale bank of
  Section~\ref{sec:frontier}, and the unbounded family is not.  The
  equivalences fix the exact missing quantifier but do not make that
  quantifier easier to prove.  \ref{res:denominator} is a finite instance of a
  \emph{different} family, the gap certificates of \ref{res:supply}, which
  reach irrationality by a one-way implication; the two families should not be
  counted together.
\item Full-depth period-multiple escape is equivalent to irrationality
  (\ref{res:fulldepth}).  The four first-harmonic pivot bounds, and the
  prime-orbit gap they refine, remain sufficient conditions and are unproved.
  Failure of the equivalent condition would settle \#249 in the opposite
  direction; failure of a sufficient route would refute only that argument.
\item Carry-rank compression and genuinely coupled denominator compression
  are separate contradiction targets.  No required upper bound or coupled
  approximant family is proved.
\item The $k$-kernel rank statement of Theorem~\ref{thm:kkernelrank} is
  independent of the irrationality problem and settles nothing about it.
  Zero-residue, composite-base arithmetic, spanning, independence, exact
  rank, basis, and relation-module dimension are Lean-checked in
  \texttt{AllBaseTotientKernel.lean}.  Martin's Theorem~1 remains prior art
  for the independence conclusion and is not formalised; it is not a remaining
  Lean gap.  The residue condition is $k\nmid r$, not $\gcd(k,r)=1$.  No source
  was located for the rank, basis, or normal form themselves; that search
  result is not a priority verdict.
\end{enumerate}

\noindent The correlation source that fits the normalised totient most
directly is Balasubramanian--Giri--Srivastav~\cite{bgs2017}, not a direct
transfer of the Tao--Ter\"av\"ainen method.  Write
\[
  g(n)=\frac{\ph(n)}n
      =\sum_{d\mid n}\frac{\mu(d)}d=(f*1)(n),
  \qquad f(d)=\frac{\mu(d)}d\in\mathcal A_1 ,
\]
the class $\mathcal A_1$ being the coefficient class in which the theorem
cited next is stated.
Theorem~2.2 of the arXiv version of~\cite{bgs2017} gives, uniformly for
$|h|\le x/2$, an explicit asymptotic for
$\sum g(n)g(n-h)$ with error $O(\log^2x)$; Remark~2.4 gives its Euler product,
For each fixed $h$, the weighted partial-summation formula immediately
following Corollary~2.8 permits $Q(n)=n(n-h)$, restoring the two linear
factors needed to return from $g(n)g(n-h)$ to
$\ph(n)\ph(n-h)$ with an explicitly propagated error.  That displayed
weighted formula is not itself stated uniformly in $h$ and starts its main
integral at \(1\), not \(H\); retaining uniformity and the lower endpoint
would require applying partial summation directly to Theorem~2.2.  No later
claim here relies on an unproved uniform weighted version.

Tao--Ter\"av\"ainen's quantitative correlation theorem
(\cite{taoteravainen2025}, Theorem~3.1) assumes either quantitative
equidistribution together with an exact small-prime condition, or
non-pretentiousness.  Here $g(p)=1-1/p$, so the stated small-prime condition
does not hold, while
$\sum_p(1-g(p))/p=\sum_p1/p^2<\infty$, so $g$ is pretentious to the constant
function.  Their theorem therefore does not apply unchanged.  The
Balasubramanian--Giri--Srivastav theorem supplies the missing first- and
second-moment input in the correct divisor-convolution class, but it does
\emph{not} supply the residue small-ball or phase anti-concentration estimate
needed to turn those moments into a certificate.  The obstacle is not the
extension of a fixed-shift theorem: Theorem~2.2 already allows $|h|\le x/2$.
The missing input is control of a nonlinear totient block whose length grows
with the truncation depth $L$, not a two-point correlation of
$\ph(n)/n$.  That modulo-$2^L$ step is the precise remaining analytic gap;
no irrationality conclusion is drawn from the correlation asymptotic alone.



% === 10_replace_status_notes_by_mathematical_boundary; original lines 1870--1916 ===
% Destination: a249_p1b.tex prop:D5cons and prop:D7-inv; a249_front.tex thm:rankonefloor and thm:cfdenom. Superseded status history is archival only.
\section*{Formal source notes}

\paragraph{Claim status.}
Two of the items above are checked propositions that no row of the claim
registry currently owns: the three declarations of \ref{res:basis}, and the
five of \ref{res:weight}; the supply implication of \ref{res:supply} is
likewise unowned, and \ref{res:denominator} cites three declarations of which
the registry's row owns one.  Under this project's own division of authority
the registry, not a manuscript, owns public status, so the honest tag on those
items is that they are kernel-checked and not yet registered.  We print that
rather than borrowing a neighbouring row's status.  The gap is a defect in the
record, not in the proofs, and closing it is a review action rather than a
mathematical one.

This manuscript is authored exposition, not proof authority.  The linked Lean
snapshot is authoritative only for its exact propositions, and kernel checking
establishes that a proposition was proved, not that it is interesting, novel,
or sufficient.  The consequence drawn in the last sentence of \ref{res:basis},
the deduction in \ref{res:periodic} that an unbounded weight is not eventually
periodic, the mechanism sentence in Section~\ref{sec:nogo}, and the size
argument in Section~\ref{sec:conditional-criteria} are one-line arguments in the prose,
marked as such; the first three are drawn from checked statements, and the
last assumes the unproved bounded-order condition.

The small numerical
instances printed above, the level-three reductions, the certificate at
$(1,12,16)$, the first values of $\Aw$, the $p=3$ Euler-factor calculation,
the mixed-difference table, and the values of $H_t$, are computations from
the definitions and checked identities beside them, and are not separate
checked propositions.  Source theorems attributed to Allouche and Shallit,
Martin, Coons, Bell and Smertnig, Kaneko--Suzuki--Tachiya, Luca--Tachiya, and
Tao--Ter\"av\"ainen are cited from the literature.  This development does not
formalise Martin's positive-density theorem or Coons's non-regularity theorem.

It does kernel-check an independent dyadic basis/rank proof and a full-kernel
infinite-dimensionality consequence that Coons already implies.  For every
integer base $k\ge2$ the same arithmetic reductions, unconditional spanning,
independence, exact rank $k^e+1$, explicit basis, and relation-module dimension
are kernel-checked in \texttt{AllBaseTotientKernel.lean}; the older conditional
module is not the present authority.  The
assessment that no prior source gives an explicit basis, an exact finite-level
dimension formula, or a relation normal form for the totient $2$-kernel is the
outcome of a literature search and is not a proof of non-existence.

Erd\H{o}s Problem~\#249 remains open.



% === 11_retain_only_cited_bibliography; original lines 1916--2000 ===
% Destination: Full bibliography remains in the long record and transferred source.
\begin{thebibliography}{10}
\bibitem{allouche-shallit} J.-P.~Allouche and J.~Shallit,
  \href{https://cs.uwaterloo.ca/~shallit/Papers/as0.pdf}{\emph{The ring of
  $k$-regular sequences}}, Theoret.\ Comput.\ Sci.\ \textbf{98} (1992),
  no.~2, 163--197,
  doi:\href{https://doi.org/10.1016/0304-3975(92)90001-V}
  {10.1016/0304-3975(92)90001-V}.  Definitions~1.1 and 2.1 are on
  pp.~2--3 of the linked author preprint, which differs slightly from the
  published version; the cited definitions are unaffected.
\bibitem{coons} M.~Coons, \emph{(Non)Automaticity of number theoretic
  functions}, J.~Th\'eor.\ Nombres Bordeaux 22 (2010), no.~2, 339--352;
  \href{https://doi.org/10.5802/jtnb.718}{doi:10.5802/jtnb.718};
  \href{https://arxiv.org/abs/0810.3709}{arXiv:0810.3709}.  Theorem~3.2,
  pp.~348--349, in the published version (Theorem~3.3, pp.~8--9, in the
  preprint): $\ph$ is not $k$-regular for any $k\ge2$.
\bibitem{martin-phi-inequalities} G.~Martin, \emph{Simultaneous inequalities
  among values of the Euler phi-function}, 2 March 2006,
  \href{https://arxiv.org/abs/math/0603053}{arXiv:math/0603053};
  \href{https://doi.org/10.48550/arXiv.math/0603053}{doi:10.48550/arXiv.math/0603053}.
  Theorem~1: for positive integers $a_1,\dots,a_k$ and integers $b_1,\dots,b_k$
  with $a_ib_j\neq a_jb_i$, and every $C>0$, the simultaneous ratio gaps
  $\ph(a_1n+b_1)/\ph(a_2n+b_2)>C,\dots$ hold on a set of positive lower density;
  Corollary~4 carries this to $\sigma$.  Cited here as a public preprint: a
  separate journal publication was not located.
\bibitem{bell-smertnig} J.~Bell and D.~Smertnig, \emph{Mahler series with
  multiplicative coefficient sequences}, 2026,
  \href{https://doi.org/10.48550/arXiv.2603.23456}{arXiv:2603.23456}.  The
  totient generating series is not $k$-Mahler for any $k\ge2$, a consequence
  of their Theorem 1.3 recorded in the introduction.
\bibitem{kanekosuzukitachiya2026} H.~Kaneko, Y.~Suzuki and Y.~Tachiya,
  \href{https://arxiv.org/abs/2601.20743}
  {\emph{Refinements of Erd\H{o}s's irrationality criterion for certain
  sparse infinite series}}, arXiv:2601.20743v1, 2026.
  Corollary~3 is on pp.~5--6 and proves irrationality for the
  \(\sigma(n)\)- and \(\ph(n)\)-in-the-exponent families; its proof is on
  pp.~18--19.
\bibitem{taoteravainen2025} T.~Tao and J.~Ter\"av\"ainen,
  \emph{Quantitative correlations and some problems on prime factors of
  consecutive integers}, arXiv:2512.01739 (submitted December 2025, revised
  April 2026).  Theorem 1.3 proves irrationality of
  $\sum_{n\ge1}\omega(n)/2^n$ at base $2$; the extension to every integer base
  and the $\Omega$ analogue are stated as remarks, with the modifications left
  to the reader.  The theorem is on p.~4 and its proof is Section~5,
  pp.~44--56, in arXiv v2.
\bibitem{bgs2017} R.~Balasubramanian, S.~Giri and P.~Srivastav,
  \emph{On correlations of certain multiplicative functions},
  J.\ Number Theory 174 (2017), 221--238,
  \href{https://doi.org/10.1016/j.jnt.2016.10.001}{DOI};
  \href{https://arxiv.org/abs/1511.02221}{arXiv:1511.02221}.
  The arXiv version's Theorem~2.2 and Remark~2.4 give the uniform shifted
  divisor-convolution correlation ($|h|\le x/2$) and Euler product used above.
  Weighted partial summation after Corollary~2.8 restores linear factors at a
  fixed $h$; it does not bound a totient block whose length grows with $L$.
\bibitem{lucatachiya2017} F.~Luca and Y.~Tachiya,
  \href{https://www.kurims.kyoto-u.ac.jp/~kyodo/kokyuroku/contents/pdf/2014-14.pdf}
  {\emph{Linear independence results for the values of divisor functions
  series}}, RIMS K\^oky\^uroku No.~2014 (2017), 138--150.
  Theorem~A on p.~139 explicitly restates their earlier Theorem~1.1 for
  nonzero purely periodic integer weights.  Theorem~1 is on p.~139, its
  examples are on p.~140, and its proof is on pp.~149--150.
\bibitem{erdos1948} P.~Erd\H{o}s,
  \href{https://users.renyi.hu/~p_erdos/1948-04.pdf}
  {\emph{On arithmetical properties of Lambert series}},
  J. Indian Math. Soc. (N.S.) \textbf{12} (1948), 63--66.
  The integer-base full-support theorem is on p.~63; the proof is on
  pp.~63--66, and the closing totient-series remark is on p.~66.
\bibitem{erdosgraham1980} P.~Erd\H{o}s and R.~L.~Graham, \emph{Old and New
  Problems and Results in Combinatorial Number Theory}, 1980, p.~61.
\bibitem{erdos1988} P.~Erd\H{o}s, \emph{On the irrationality of certain
  series: problems and results}, in A.~Baker (ed.), \emph{New Advances in
  Transcendence Theory}, Cambridge UP, 1988, pp.~102--109,
  doi:\href{https://doi.org/10.1017/CBO9780511897184.009}
  {10.1017/CBO9780511897184.009}.
\bibitem{fan2026totient} S.~Fan, comment on Erd\H{o}s Problem \#249,
  \href{https://www.erdosproblems.com/forum/thread/249}{\texttt{erdosproblems.com/forum/thread/249}},
  19:01 on 16 May 2026: the convolution identity $\varphi=\mu*\mathrm{id}$
  yields $\sum_{n\ge1}\varphi(n)/2^n=\sum_{n\ge1}\mu(n)/(2^n-1)^2+1/2$.
\bibitem{erdosproblems} T.~F. Bloom,
  \href{https://www.erdosproblems.com/249}{\emph{Erd\H{o}s Problem \#249}},
  \texttt{erdosproblems.com/249}, accessed 28 July 2026 (page displays
  ``last edited 28 September 2025'').  The current record labels the problem
  open, cites \texttt{[ErGr80, p.~61]} and \texttt{[Er88c, p.~102]}, and
  explicitly describes its status as the website owner's present assessment
  rather than a literature-completeness guarantee.
\end{thebibliography}
